% =====================================================================
%  main_paper_ieeetj.tex -- IEEE Transactions (ieeetj.cls) submission version.
%  Converted from main_paper.tex (article.cls working draft) using the official
%  IEEE Transactions template (ieeetj.cls, August 2022) supplied with this repo.
%  Compile:  cd theory && pdflatex main_paper_ieeetj.tex   (three passes)
%  NOTE: fill in received/revised/accepted/published/currentdate/doiinfo and
%        supply author photos before final submission (see conversion notes).
% =====================================================================
\documentclass{ieeetj}
\usepackage{amsmath,amssymb,amsthm}
\usepackage{booktabs}
\usepackage{pgfplots}
\usepackage{pgfplotstable}
\usepackage[T1]{fontenc}
\usepackage[utf8]{inputenc}
\usepackage{hyperref}
\hypersetup{hidelinks}

%  Colors and header logos required by ieeetj.cls (from the template preamble)
\AtBeginDocument{\definecolor{tmlcncolor}{cmyk}{0.93,0.59,0.15,0.02}%
\definecolor{NavyBlue}{RGB}{0,86,125}}
\def\OJlogo{\vspace{-4pt}}
\def\seclogo{\vspace{-4pt}}
\def\authorrefmark#1{\ensuremath{^{\textbf{#1}}}}

%  Manuscript-level macros (kept identical to the article.cls draft)
\providecommand{\RDP}{\mathrm{RDP}}
\def\figdatadir{../figures/dat/}
\pgfplotsset{compat=1.18}

\theoremstyle{plain}
\newtheorem{theorem}{Theorem}
\newtheorem{proposition}{Proposition}
\newtheorem{lemma}{Lemma}
\newtheorem{corollary}{Corollary}
\theoremstyle{definition}
\newtheorem{assumption}{Assumption}
\newtheorem{remark}{Remark}

\begin{document}
\receiveddate{XX Month, XXXX}
\reviseddate{XX Month, XXXX}
\accepteddate{XX Month, XXXX}
\publisheddate{XX Month, XXXX}
\currentdate{XX Month, XXXX}
\doiinfo{XXXX.XXXX.XXXXXXX}
\markboth{}{Q. Wang {et al.}}

\title{Auditable Privacy Accounting and Dynamic Semantic Budget
Allocation for Differentially Private Semantic Communications}

\author{Qinmin Wang\authorrefmark{1}, Zixin Wang\authorrefmark{2},
Jie Peng\authorrefmark{3}, Fenggui Pei\authorrefmark{1},
and Zongzhong Wang\authorrefmark{4}}
\affil{School of Computer Science and Artificial Intelligence, Guangdong
University of Science and Technology, Dongguan 523088, China}
\affil{College of Electronic and Communication Engineering, Tianjin Normal
University, Tianjin 300387, China}
\affil{Shenzhen Customs, Shenzhen 518000, China}
\affil{School of Electromechanical Engineering, Guangdong University of
Technology, Guangzhou 510006, China}
\corresp{Corresponding author: Zongzhong Wang (email: wangzz@gdut.edu.cn).}
\authornote{This work was supported in part by the 2025 Science and Technology
Innovation Capability Enhancement Project of Guangdong University of Science
and Technology under Grant No.~GKY-2025BSQDK-46. E-mail addresses of the other
authors: wangqinmin@gdust.edu.cn, fengguipei@gmail.com, 396514839@qq.com,
35743222@qq.com. ORCID: 0009-0006-0447-0482 (Q. Wang), 0009-0006-7807-3666
(Z. Wang).}

\maketitle

\begin{abstract}
Decomposed latents in semantic communication retain a subject's identity, yet
existing designs certify protection from the \emph{appearance} of the released
noise rather than from an auditable budget, and report no adversary
identification bound. We develop a framework that quantifies both. The
per-component $\ell_2$ sensitivity is isolated and correct quadrature
bookkeeping is distinguished from looser linear practices: at a low-noise
operating point, amplitude-only $\ell_1$ accounting reports
$\varepsilon\approx693$ where the rigorous Gaussian tail bound gives
$\approx2441$ and the quadrature account gives $2000$ (a Gaussian-proxy
re-pricing of a reference operating point). Two independent guarantees are
derived. A $(\varepsilon,\delta)$-DP account via R\'enyi and Gaussian DP
audits the release; a Gaussian worst-case mutual-information bound with Fano's
inequality gives an identification-error floor $\tau$ with ceiling
$\tau_{\max}(K)$. Neither implies the other. Post-processing leaves both
unchanged, so a public invertible shaping network cannot help Bob without
helping Eve. The induced budget program has a closed-form water-filling KKT
solution. On a SemanticStyleGAN calibration, $65{,}536$ prior codes fix each
block's power and sensitivity and noise sweeps fix utility weights; the
accounting ordering, its dependence on task--private overlap, and the failure
of amplitude-only accounting are reproduced. The perceptual score and attack
error remain surrogates; source statistics are a pushforward of the generator
prior, not an inversion of real identities.
\end{abstract}

\begin{IEEEkeywords}
Semantic communications, differential privacy, R\'enyi differential privacy,
Gaussian differential privacy, mutual information, Fano inequality, resource
allocation, water-filling, latent representations.
\end{IEEEkeywords}

% =====================================================================
\section{INTRODUCTION}
\label{sec:intro}
% =====================================================================

\IEEEPARstart{S}{emantic} communication compresses the source to a latent
representation that carries the task-relevant meaning rather than the raw
waveform. When the source is a face, that representation retains identity:
the encoder separates global appearance and local part codes precisely so that
the image can be reconstructed or edited, and the same part codes an eavesdropper
can observe from the release. A privacy mechanism that adds noise to a subset
of these codes therefore has to answer two questions that are usually left
implicit: what exactly is being certified, and how much legitimate distortion
the certificate costs.

The scheme of~\cite{arXiv2510} perturbs a subset of semantic latent codes, then uses a discriminator
to shape that perturbation. The goal is twofold: keep the released noise close to a
Laplace draw, and keep the distortion semantically mild. It is a reasonable engineering
trade-off. That work \emph{does} state a nominal privacy budget --- a global sensitivity $\Delta f=351.88$,
with Laplace noise $\mathrm{Lap}(0,\Delta f/\varepsilon)$ calibrated for $\varepsilon\in\{1,5,\ldots,2000\}$. What is missing is not a parameter but an
\emph{auditable guarantee for the noise actually released}. The discriminator learns a pattern that only approximates a Laplace draw; the
published $(\varepsilon,\delta)$ value is therefore not a proof about the released mechanism, and no
identification-error statement is made. This work does not replace that design. It
supplies a framework to audit any such release under standard DP metrics and to check it
against an identification attack. For comparison we re-price a \emph{Gaussian proxy} of its operating
point (Section~\ref{sec:experiments}): how the budget should be charged if the pattern were an exact
additive Gaussian mechanism. This is a statement about the accounting law, not about the
pattern's true distribution.

The framework reports two complementary evaluation tools that must not be
confused. The $(\varepsilon,\delta)$-DP system (RDP, GDP, linear tail bound)
guarantees that adjacent records are indistinguishable. The Fano certificate
gives an operational lower bound on an adversary's identification error,
derived from second-order release statistics. One does not imply the other: a
release can satisfy a DP budget yet still leak an identifiable amount, and a
release can hold a Fano floor without satisfying DP. The two are reported side
by side throughout.

This gap is not unique to the reference scheme. Most existing designs optimize a learned
noise-shaping network against a distortion objective; they rarely state the
component-wise sensitivity of the decomposed latent, fix a composition law for
the private blocks, or report a DP budget and an identification-error floor
jointly. The present work starts from the sensitivity and the composition law
so that the two guarantees can be priced and compared on the same release.

The contributions fall into four parts.
\emph{(1) Sensitivity and bookkeeping.} The component-wise $\ell_2$ sensitivity
of a decomposed latent is isolated, and the three bookkeeping rules that follow
separate a correct per-component quadrature account from the looser monolithic
and per-component-global practices; this clarifies how a privacy budget is
charged to individual semantic components.
\emph{(2) Post-processing and invertible shaping.} An invertible learned shaping
cannot raise the guarantee: post-processing leaves the privacy loss unchanged,
so such a component is legitimate only as a keyed, key-conditional design
element. This answers whether a learned invertible denoiser can improve Bob's
utility without also helping Eve. The certificate below assumes $T$ invertible
and known to the eavesdropper (the worst case); a Bob-secret $T$ is reported as
key-conditional rather than plain $(\varepsilon,\delta)$-DP.
\emph{(3) Identification certificate.} A Gaussian worst-case mutual-information
bound combined with Fano's inequality certifies an identification-error floor
for an adversary observing the release, with an explicit feasibility ceiling
$\tau_{\max}(K)$. This supplies an operational attack-error lower bound in
place of a marginal-distribution ``looks private'' argument.
\emph{(4) Allocation law and DP audit.} The certificate's budget program is
convex and admits a closed-form water-filling KKT solution with two asymptotic
regimes; the same KKT analysis gives the noise profile induced by each standard
accounting law, and the static-index policy is a degenerate corner of the
program. The standard R\'enyi-DP to $(\varepsilon,\delta)$ conversion and the
Gaussian-DP quadrature composition are derived, and the classical linear tail
bound is audited against them.

\emph{Scope of the evaluation.} The sensitivity derivation, post-processing theorem, Fano certificate, KKT
allocation and composition rules are mechanism-level results. They hold for any release
that satisfies the model assumptions, not just for the generator used below. The
quantitative calibration does rest on the released SemanticStyleGAN generator: per-block
source power, sensitivity, utility weight and shaping residual are measured on real
decomposed codes and generator decodes (Section~\ref{sec:latent}). Three findings are reproduced on
that measured calibration --- the accounting ordering, the dependence on the
task--private overlap, and the failure mode of the amplitude-only law. Two quantities
remain documented surrogates. The LPIPS score is a model-based prediction on generator
decodes; the identification error against the Fano floor awaits an inverted-identity
corpus with labels, as recorded in Section~\ref{sec:discussion}.

The remainder of the paper is organized as follows.
Section~\ref{sec:related} positions the work against semantic-privacy, DP
accounting, and information-theoretic identification literature.
Section~\ref{sec:setup} fixes the latent decomposition, adjacency, and notation.
Section~\ref{sec:sensitivity} derives the component sensitivity and the
post-processing result. Section~\ref{sec:certificate} states the Fano
identification certificate. Section~\ref{sec:allocation} solves the budget
allocation, and Section~\ref{sec:rdpgdp} derives the RDP/GDP composition and the
linear tail audit. Section~\ref{sec:experiments} reports the anchor,
image-proxy and latent-level evaluations, and Section~\ref{sec:discussion}
collects the limitations and future work.

The evaluation uses three calibrations. The first is anchored to the reference
backbone and exercises every accounting rule end to end; the second is built
from 500 CelebAMask-HQ faces~\cite{celebamaskhq}, with per-part variance and
clipped range measured on real masks; the third is a latent-level calibration
on the released SemanticStyleGAN generator, where all four calibration
scalars (source power, sensitivity, utility weight and shaping residual) are
measured on real decomposed codes and generator decodes. On all three, the
price of a certificate is governed by the task--private overlap rather than by
the reported $\varepsilon$, the amplitude-only $\ell_1$ law under-reports
$\varepsilon$ at the low-noise (large-$\mu$) operating point, and a
static index selection certifies no identification floor. The one remaining
surrogate is the identification attack error, whose measurement against the
certified floor requires an inverted-identity corpus; the stated limitations in
Section~\ref{sec:discussion} record this explicitly.

% =====================================================================
\section{RELATED WORK}
\label{sec:related}
% =====================================================================

\emph{Semantic representations and their privacy.} A modern generative encoder splits an image into a global structure code and
part-wise local codes, which is what makes controlled synthesis and editing
possible~\cite{semanticstylegan}. Over a wireless link, releasing those codes is the standard
semantic-communication step --- and the part codes carry the subject's identity. Two
surveys map the threat models and defense techniques in this setting~\cite{semcomsurvey_won,semcomsurvey_meng}. The reference
design~\cite{arXiv2510} perturbs a subset of codes with Laplace noise, calibrated to a stated global
sensitivity ($\Delta f=351.88$) and nominal $\varepsilon$, then shapes the result with a learned pattern to hold
down Bob's distortion. Two things remain implicit: the \emph{component-wise} sensitivity of the decomposed
code, and the composition law across private blocks. Nor does the learned released noise
carry an auditable DP proof. The same pattern recurs in much semantic-privacy work,
where a learned noise shape is tuned against reconstruction distortion (e.g.\ LPIPS)
with no identification attack or error floor evaluated. \emph{Gap:} what is missing is an
auditable, block-wise account paired with an operational attack bound --- not the mere
presence of an $\varepsilon$ value. This paper starts from the component sensitivity and the
composition law.

\emph{Privacy accounting.} The second reference point is the DP accounting
literature. The Gaussian mechanism and its classical linear
composition are standard~\cite{dwork2014}. R\'enyi DP accumulates the
R\'enyi divergence of the Gaussian outputs quadratically and converts to
$(\varepsilon,\delta)$ with an order-optimised bound~\cite{mironov2017rdp}.
Gaussian DP treats the privacy-loss random variable as Gaussian and composes
its mean and variance in quadrature~\cite{dong2022gdp}. These tools were
developed and systematised alongside differentially private
training~\cite{abadi2016}, where the queried statistic is a monolithic model
gradient. They do not transfer unchanged to a decomposed semantic latent: the
sensitivity must be charged per block, the private subset is a design choice,
and a learned invertible post-processing sits between the release and the
receiver. The three bounds are applied per semantic component rather than to a
monolithic vector, and compared on the same release. Membership inference
attacks~\cite{shokri2017} ask whether a particular record contributed to the
release; the identification certificate here is a different $K$-ary problem,
namely which enrolled template the released code matches.

\emph{Information-theoretic privacy and identification.} Fano's inequality offers a third reference point: it converts an information bound
into a lower bound on the error of a $K$-ary estimator~\cite{fano}. Information-bottleneck and
privacy-funnel formulations~\cite{privacyfunnel} trade released information against task utility by
optimizing an objective. This work does not optimize such a loss. It instead derives an
auditable, verifiable lower bound on an adversary's identification error from
second-order release statistics. The bound is stated in operational terms the system
cares about: the error an adversary cannot drive below, rather than a marginal
distribution that merely looks private.

% =====================================================================
\section{SYSTEM MODEL AND NOTATION}
\label{sec:setup}
% =====================================================================

\emph{Semantic latent.} A semantic encoder (here a StyleGAN generator inverted
by a Semantic StyleGAN encoder~\cite{semanticstylegan}) maps an image to a
latent $z\in\mathbb{R}^{n}$ decomposed into $H$ semantic components
\begin{equation}
 z=(z_1,\dots,z_H),\qquad z_h\in\mathbb{R}^{d_h},\qquad
 n=\textstyle\sum_{h=1}^{H} d_h ,
\label{eq:decomp}
\end{equation}
where $h$ indexes a shared (global) code or a local (part-wise) code group.
For the reference backbone, $H_{\mathrm c}=28$ codes of dimension
$d_h=512$, $n=14336$; our implementation groups the 13 part-wise local codes
into 13 blocks of $d_h=1024$ plus 2 shared blocks of $d_h=512$ ($H=15$,
$n=14336$). Grouping affects only bookkeeping; the composition formulas of
Section~\ref{sec:rdpgdp} are additive per component and are unchanged by it.

\begin{table}[h]
\centering\footnotesize
\setlength{\tabcolsep}{4pt}
\caption{Core notation.}
\label{tab:notation}
\begin{tabular}{cp{4.4cm}}
\toprule
symbol & meaning \\
\midrule
$z_h,\,d_h$ & block latent vector and its width \\
$g_h$ & block source power $\operatorname{Tr}\Sigma_h$ \\
$\Delta f_h$ & block $\ell_2$ sensitivity \\
$\sigma_h$ & released noise standard deviation \\
$P,\,Q$ & private block set and its complement \\
$H_{\mathrm c},\,H$ & original code count $28$, grouped block count $15$ \\
$K,\,\tau$ & gallery size and certified error floor \\
$\nu$ & KKT dual multiplier; marginal cost per unit leakage \\
$\mu_h$ & per-block privacy-loss parameter $\Delta f_h/\sigma_h$ \\
$B$ & private identity template, one of the $K$ enrolled templates \\
$a_h=w_h\kappa_h$ & block task-cost coefficient in the allocation \\
$\Sigma_h$ & block latent covariance matrix \\
$c$ & linear tail constant $\sqrt{2\ln(1.25/\delta)}$ \\
$S$ & static protected block set of size $k$ \\
\bottomrule
\end{tabular}
\end{table}
Table~\ref{tab:notation} collects these symbols for reference. Throughout,
$n$ denotes the total latent coordinate count and $\ln,\log_2$ denote natural
and base-two logarithms; mutual information is reported in bits and the
allocation program accumulates nats (Remark~\ref{rm:units}).

\emph{Private subspace.} Let $P\subseteq\{1,\dots,H\}$ be the set of
\emph{declared private} components and $Q=\{1,\dots,H\}\setminus P$ its
complement. $P$ is an input design decision made by the system operator, not a
quantity learned or output by the allocation algorithm: it lists the semantic
attributes whose inference the system promises to obstruct (identity, in the
reference task). We write $z_P=(z_h)_{h\in P}$.

\begin{assumption}[Bounded coordinate range]
\label{as:range}
Every latent coordinate satisfies $z_{h,i}\in[a_h,b_h]$; without loss of
generality $b_h-a_h=b-a$ for all $(h,i)$, after the $1$-Lipschitz
normalization applied by the encoder. This common-range normalization is a
modeling assumption rather than a property measured on the reference encoder
weights. If a real encoder leaves coordinate ranges heterogeneous, replacing
the uniform $b-a$ by the measured per-coordinate range can only tighten the
sensitivity when the true ranges are smaller; the uniform value is retained
here so the sensitivity stays a safe upper bound, at the cost of possibly
over-estimating $\Delta f_h$.
\end{assumption}

\begin{assumption}[Block-wise second-order structure]
\label{as:indep}
For each block $h$, the coordinate variances are uniformly bounded,
$\operatorname{Var}(z_{h,i})\le g_{h,i}$ for every coordinate $i$, and the
block source power satisfies
$\operatorname{Tr}\Sigma_h=\sum_i\operatorname{Var}(z_{h,i})\le g_h$. No PSD
ordering $\Sigma_h\preceq\operatorname{diag}(g_{h,i})$ is assumed; residual
within-block and across-block correlation is charged conservatively to the
block source power $g_h$, see Remark~\ref{rm:corr}. Correlated
coordinates share information, which lowers the true joint mutual information;
the bound used here ignores this reduction and reports a safe upper bound on
the leakage. Correspondingly, the certified floor $\tau$ is a lower bound on
the identification error that a tighter, joint-covariance certificate could
establish: it is a safe-side certification, but it may overstate the noise cost
and make the utility look pessimistic. The relaxation never weakens the privacy
guarantee.
\end{assumption}

\emph{Vector-release convention.} Each block $h$ is released as a
$d_h$-dimensional vector. The Gaussian worst-case mutual information uses the
AM--GM (Jensen) bound
$I(z_h;z_h+N_h)\le\tfrac{d_h}{2}\log(1+\operatorname{Tr}\Sigma_h/(d_h\sigma_h^2))$
(Lemma~\ref{lm:maxent}), so $g_h$ denotes the block source power
$\operatorname{Tr}\Sigma_h$; the coordinate count $d_h$ enters both this leak
bound and the block $\ell_2$ sensitivity $\Delta f_h=\sqrt{d_h}(b-a)$.

\emph{Adjacency and code-level change.} A record is one face image and its
associated latent vector. Two records are adjacent, $D\sim D'$, if they differ
in the private semantics of exactly one subject (replace-one); every other
subject's code is held fixed. Formally, for records $D=(z^{(1)},\ldots,z^{(m)})$
and $D'$ differing in subject $j$,
\begin{equation}
D\sim D'\iff z^{(i)}=z'^{(i)}\ \forall i\ne j
\quad\text{and}\quad z_P^{(j)}\ne z_P'^{(j)},
\label{eq:adjacency}
\end{equation}
i.e.\ adjacency replaces the identity-bearing private semantics of one subject.
This differs from classical database DP, where adjacency flips one data
\emph{record}; here the unit of change is one subject's identity code inside an
otherwise fixed release, and every other subject's code is held fixed.
The reference latent has $H_{\mathrm c}=28$ codes of dimension $512$, grouped
into the $H=15$ blocks above; a local block concatenates two such codes. The
replace-one adjacency changes only the codes that carry the subject's identity;
if a concatenated block contains one changed code and one unchanged code, its
true $\ell_2$ sensitivity is $(b-a)\sqrt{512}$, not $(b-a)\sqrt{1024}$. The
block sensitivity used below takes the conservative worst case in which every
coordinate of the block changes under adjacency,
$\Delta f_h=(b-a)\sqrt{d_h}$; this over-estimates the sensitivity of a
partially changed block by at most $\sqrt{2}$, and the resulting privacy
accounting remains an upper bound. All quantities are per-record.

\emph{Release mechanism.} The transmitter releases
\begin{equation}
Y=T\!\left(z+N\right),\qquad
N_h\sim\mathcal{N}\!\left(0,\sigma_h^{2} I_{d_h}\right)\ \text{independent},
\label{eq:mechanism}
\end{equation}
where $T$ is an optional transformation. The baseline scheme
of~\cite{arXiv2510} instantiates $T$ as a learned, per-component shaping
network trained against a discriminator, with the additional requirement that
$T$ be invertible at the legitimate receiver. Two deployment cases must be
distinguished. \emph{Case (i), worst case:} $T$ is invertible and known to the
eavesdropper, so the adversary observes the same release $Y$ as Bob; this is the
standard $(\varepsilon,\delta)$-DP setting and the case certified below.
\emph{Case (ii), keyed receiver:} $T$ is a Bob-secret key unknown to the
eavesdropper. Post-processing still cannot raise the guarantee, but any
privacy statement valid only under the hidden key is \emph{key-conditional}
privacy; under the Kerckhoffs principle it is \textbf{not plain
$(\varepsilon,\delta)$-DP}, and it is reported as such.

\emph{Adversary model.} The certificate and audit below assume the adversary
observes one release $Y$ per record. In Case~(i) the adversary knows $T$ and
can apply $T^{-1}$; it may hold side information about the enrolled gallery,
and the Fano bound remains valid for any estimator that conditions on such
side information, through data processing. The account is per-record and does
not compose over repeated queries to the same subject; a multi-release adversary
would require the composition rules of Section~\ref{sec:rdpgdp} applied over
those releases, which is outside the single-release analysis here.

\begin{table}[h]
\centering\footnotesize
\setlength{\tabcolsep}{3pt}
\caption{Two deployment modes for the shaping transform $T$.}
\label{tab:modes}
\begin{tabular}{p{2.2cm}p{2.1cm}p{2.4cm}}
\toprule
mode & Eve knows $T$? & guarantee type \\
\midrule
Case (i): public/invertible & yes & standard $(\varepsilon,\delta)$-DP;
certified below \\
Case (ii): Bob-secret key & no & key-conditional privacy; \emph{not} plain DP
under Kerckhoffs \\
\bottomrule
\end{tabular}
\end{table}
Table~\ref{tab:modes} summarises the two modes; only Case~(i) is certified in
this paper, and a keyed design must be reported as key-conditional rather than
as a standard DP guarantee.

% =====================================================================
\section{SENSITIVITY AND INVERTIBLE SHAPING}
\label{sec:sensitivity}
% =====================================================================

\begin{proposition}[Accounting, not mechanism, inflates the noise]
\label{pr:inflation}
Under Assumption~\ref{as:range}, the L2 sensitivity of the whole latent and of
a single component are
\begin{equation}
\Delta f_{\mathrm{glob}}=(b-a)\sqrt{{\textstyle\sum_h d_h}},
\qquad
\Delta f_h=(b-a)\sqrt{d_h}.
\label{eq:sens}
\end{equation}
The bookkeeping ratios below are stated on the original code count
$H_{\mathrm c}=28$; after grouping into the $H=15$ blocks of
Section~\ref{sec:setup}, the additive composition formulas of
Section~\ref{sec:rdpgdp} are unchanged. Two composition concepts must be kept
separate. An \emph{amplitude-sum} account charges the sum of the per-block
sensitivity-to-noise ratios, $\sum_h\Delta f_h/\sigma_h$, which is an $\ell_1$
combination across blocks. A \emph{quadrature} account accumulates the squared
ratios, $\sum_h(\Delta f_h/\sigma_h)^2$, which for equal noise equals one global
$\ell_2$ sensitivity query; this matches the Gaussian RDP/GDP composition of
Section~\ref{sec:rdpgdp}. Three bookkeeping rules follow. Rules (i) and (iii)
are the loose (and, at large $\mu$, unsound) practices that this paper audits;
their $5.29\times$ ratios below are the gap between incorrect and correct
bookkeeping, not a gain of any noise mechanism. In all three the symbol
$\varepsilon$ denotes a noise-to-sensitivity ratio; the DP parameter further
requires the tail constant $c=\sqrt{2\ln(1.25/\delta)}$ of
\eqref{eq:tailconst}, suppressed here for a clean ratio comparison.
\begin{enumerate}
\item[(i)] \emph{per-component with $\ell_1$ composition:}
$\varepsilon=\sum_{h}\Delta f_h/\sigma_h$, so the smallest equalising noise
satisfies $\sigma\ge(b-a)\textstyle\sum_h\sqrt{d_h}/\varepsilon$;
\item[(ii)] \emph{monolithic vector query:}
$\sigma\ge\Delta f_{\mathrm{glob}}/\varepsilon
=(b-a)\sqrt{{\textstyle\sum_h d_h}}/\varepsilon$;
\item[(iii)] \emph{charging every component the global sensitivity} (the
practice implied when one reported $\varepsilon$ is reused for every code):
$\sigma\ge H_{\mathrm c}\,\Delta f_{\mathrm{glob}}/\varepsilon$.
\end{enumerate}
For $d_h\equiv d$, the ratios are
$\sigma_{\mathrm{(i)}}/\sigma_{\mathrm{(ii)}}=\sqrt{H_{\mathrm c}}$ and
$\sigma_{\mathrm{(iii)}}/\sigma_{\mathrm{(i)}}=\sqrt{H_{\mathrm c}}$; the
third-to-first factor is $H_{\mathrm c}$ in noise power, or
$10\log_{10}H_{\mathrm c}$ decibel.
\end{proposition}

\begin{proof}
(i) is the standard Gaussian-mechanism bound applied blockwise, using
$\|\Delta z_h\|_2\le(b-a)\sqrt{d_h}$ on every adjacent pair.
(ii) follows from $\|\Delta z\|_2^2=\sum_h\|\Delta z_h\|_2^2
\le(b-a)^2\sum_h d_h$. (iii) substitutes $\Delta f_{\mathrm{glob}}$ for every
summand in (i). For $d_h\equiv d$:
$\sum_h\sqrt d=H_{\mathrm c}\sqrt d$ and
$\sqrt{\sum_h d}=\sqrt{H_{\mathrm c}d}$, giving
$\sigma_{\mathrm{(i)}}/\sigma_{\mathrm{(ii)}}=\sqrt{H_{\mathrm c}}$, and
$\sigma_{\mathrm{(iii)}}/\sigma_{\mathrm{(i)}}=
H_{\mathrm c}\sqrt{H_{\mathrm c}d}/(H_{\mathrm c}\sqrt d)=\sqrt{H_{\mathrm c}}$.
\end{proof}

\emph{Illustrative numbers.} Take $H_{\mathrm c}=28$ codes of width
$d=512$, range $b-a=1$, and a unit budget $\varepsilon=1$ with equal noise
across blocks. The three rules give the required common standard deviation
$\sigma_{\mathrm{(i)}}=28\sqrt{512}\approx634$,
$\sigma_{\mathrm{(ii)}}=\sqrt{28\cdot512}\approx120$, and
$\sigma_{\mathrm{(iii)}}=28\sqrt{28\cdot512}\approx3352$. The gap between
(i) and (ii) is purely a composition choice: rule (i) adds the 28 block
amplitudes linearly, $\sum_h\sqrt d=28\sqrt d$, whereas rule (ii) combines the
block sensitivities in quadrature, $\sqrt{\sum_h d}=\sqrt{28d}$; the square-root
of a sum is smaller than the sum of square roots, which is why (i) demands about
$5.29\times$ more noise amplitude than (ii). Rule (iii) asks for another
$5.29\times$ on top of (i), and $28\times$ (14.5\,dB) the noise power of (i).
Rule (ii) already matches the correct quadrature account.

\textbf{The gain exposed by Proposition~\ref{pr:inflation} is a correction of
bookkeeping, not a new noise mechanism: a correct monolithic vector account
(ii) is as tight as correct per-component quadrature accounting
(Proposition~\ref{pr:quadrature}). Any larger privacy--utility dividend must
come from the allocation itself (Section~\ref{sec:allocation}), which is a
separate degree of freedom.}

\begin{theorem}[Post-processing and invertible shaping]
\label{th:postproc}
Fix the noise law in \eqref{eq:mechanism}.
\begin{enumerate}
\item[(i)] For every $T$ that is a (possibly randomized) function of the
mechanism output with randomness independent of $D$, $T\circ Y$ is
$(\varepsilon,\delta)$-DP whenever $Y$ is.
\item[(ii)] If $T$ is invertible and $T^{-1}$ is available to the adversary,
$T\circ Y$ provides exactly the same guarantee as $Y$: the adversary recovers
$Y$ and the protection, if any, is due to $N$ alone.
\item[(iii)] If $T$ depends on the private latent directly rather than on the
release, it is part of the mechanism and its sensitivity must be re-accounted;
$T$ may then \emph{increase} $\varepsilon$.
\end{enumerate}
\end{theorem}

\begin{proof}
(i) is the post-processing inequality: for adjacent $D\sim D'$ there is a
coupling $(Y,Y')$ of the two output laws with $\Pr[Y\neq Y']\le\delta$;
applying the same $T$ to both sides preserves
$\{T(Y)=T(Y')\}\supseteq\{Y=Y'\}$, hence
$\Pr[T(Y)\neq T(Y')]\le\delta$, and on the complement the two conditional
laws coincide, which is exactly $(\varepsilon,\delta)$-DP.
(ii) Eve computes $T^{-1}$, so the information available to her is the
$\sigma$-field generated by $Y$; by the data-processing inequality
$I(z;T^{-1}(T(Y)))=I(z;Y)$ and every distinguishing advantage is unchanged.
(iii) $T(z+\cdot)$ is a function of the data, so the pair
$(z,T(z+N))$ is released by a mechanism whose sensitivity is governed by
$\|T(z)-T(z')\|$, not by $\Delta f$. A per-image contrast normalization
$\sigma\mapsto\sigma/\widehat{s}(z)$ is an instance: the released scale factor
$1/\widehat{s}(z)$ is computed from the record, so it leaks through
$\widehat{s}(z)$. If $\widehat{s}(z)\in[s_{\min},s_{\max}]$ with $s_{\min}>0$,
its scalar sensitivity is
$|1/\widehat{s}(z)-1/\widehat{s}(z')|\le |\widehat{s}(z)-\widehat{s}(z')|/
s_{\min}^{2}$; this quantity must be added to $\Delta f_h$ and re-accounted,
and the normalization cannot be treated as a fixed, privacy-neutral constant.
If $s_{\min}=0$, the gain $1/\widehat{s}(z)$ is unbounded and no finite DP
guarantee holds for that release.
\end{proof}

\emph{How to re-account such a scaling.} If a record-dependent gain
$g(z)=1/\widehat{s}(z)$ is released alongside the block, it is treated as an
additional scalar output with sensitivity
$\Delta g=|\widehat{s}(z)-\widehat{s}(z')|/s_{\min}^{2}$; this $\Delta g$ is
combined with the block sensitivity $\Delta f_h$ under the chosen composition
rule (linearly for the amplitude-sum account, in quadrature for the RDP/GDP
account), and the allocation is re-solved with the enlarged sensitivity. The
mechanism and calibration in this paper use a record-independent per-block
standard deviation $\sigma_h$; no per-image contrast or gain normalization is
applied, so this extra sensitivity term is zero throughout the evaluation.

Part~(ii) has a direct operational consequence: a learned
invertible network cannot denoise a differentially private release for Bob
alone. If Bob's decoder removes the residual noise well enough to recover a
useful signal, the same invertible map applied by Eve recovers the same
residual; Bob's advantage over Eve cannot come from denoising a DP release. The
only way to evade this conclusion is a shaping $T$ that Eve does not know, and
that key-conditional design is not standard $(\varepsilon,\delta)$-DP. The
remaining origin of a legitimate advantage is noise whose projection onto the
task subspace vanishes, which is the allocation problem of
Section~\ref{sec:allocation}.

\begin{remark}[The invertibility--privacy tension]
\label{rm:tension}
Part~(ii) imposes a constraint that applies to any deployment: DP noise cannot
be ``learned away'' for the legitimate receiver only. The deprotection network
at Bob is a decoder of the released signal; if it recovers $z$ up to a residual
that is useful to Bob, the same architecture applied by Eve recovers the same
residual. Hence Bob's advantage over Eve cannot originate in denoising a
differentially private release. It can only originate in
(a) a \emph{secret} shaping $T$ that Eve does not know, in which case the
guarantee is key-conditional and must be stated as such; it is not
$(\varepsilon,\delta)$-DP in the Kerckhoffs sense; or
(b) noise whose projection onto the task subspace is zero, in which case Bob
needs no denoiser and the protocol reduces to an allocation over the
components that actually project onto the task subspace.
The reference scheme~\cite{arXiv2510} relies on a learned noise pattern to
keep Bob's distortion low at large $\varepsilon$. Under
Theorem~\ref{th:postproc} and this remark, if the pattern is public its privacy
contribution can only come from the added noise itself; a low-distortion release
attributed to a public learned denoiser should be re-accounted under the noise
alone, and a Bob-secret pattern should be reported explicitly as key-conditional
rather than as plain $(\varepsilon,\delta)$-DP.
\end{remark}

% =====================================================================
\section{A CERTIFICATE FOR THE PRIVATE SUBSPACE}
\label{sec:certificate}
% =====================================================================

Data processing bounds the leakage of any decision Eve runs on the release.
This bound converts a privacy budget into an operational error bound for the
attack family the system cares about: identification among $K$ enrolled
templates.

\begin{lemma}[Gaussian worst case at fixed covariance]
\label{lm:maxent}
Let $z_h\in\mathbb{R}^{d_h}$ be the released block with covariance $\Sigma_h$
and block source power $g_h=\operatorname{Tr}\Sigma_h$, and let
$N_h\sim\mathcal{N}(0,\sigma_h^{2}I_{d_h})$ be independent of $z_h$. Then
\begin{align}
I(z_h;z_h+N_h)
&\le\tfrac12\log\det\!\left(I+\frac{\Sigma_h}{\sigma_h^{2}}\right)\notag\\
&\le\tfrac{d_h}{2}\log\!\left(1+\frac{g_h}{d_h\,\sigma_h^{2}}\right).
\label{eq:maxent}
\end{align}
\end{lemma}

\begin{proof}
Write $Y_h=z_h+N_h$. Among random vectors with fixed covariance $\Sigma_h$, the
Gaussian achieves the maximal differential entropy, so
$h(z_h)\le\tfrac12\log\big((2\pi e)^{d_h}\det\Sigma_h\big)$ and
$h(Y_h)\le\tfrac12\log\big((2\pi e)^{d_h}\det(\Sigma_h+\sigma_h^{2}I)\big)$;
the conditional entropy is
$h(Y_h\mid z_h)=h(N_h)=\tfrac12\log\big((2\pi e\sigma_h^{2})^{d_h}\big)$. Hence
\begin{equation*}
I(z_h;Y_h)=h(Y_h)-h(Y_h\mid z_h)
\le \tfrac12\log\det\!\left(I+\frac{\Sigma_h}{\sigma_h^{2}}\right).
\end{equation*}
Let $\lambda_i\ge0$ be the eigenvalues of $A=\Sigma_h/\sigma_h^{2}$. By the
arithmetic--geometric mean inequality,
$\det(I+A)=\prod_{i=1}^{d_h}(1+\lambda_i)
\le\big(\tfrac{1}{d_h}\sum_{i=1}^{d_h}(1+\lambda_i)\big)^{d_h}
=\big(1+\tfrac{1}{d_h}\operatorname{Tr}A\big)^{d_h}$, which yields the second
inequality in \eqref{eq:maxent}. (The commonly quoted $\det(I+A)\le1+\operatorname{Tr}A$
is \emph{false} for $d_h>1$: for $A=I_2$, $\det(I+A)=4>3=1+\operatorname{Tr}A$;
in fact $\det(I+A)\ge1+\operatorname{Tr}A$, so $1+\operatorname{Tr}A$ cannot serve
as an upper bound.) The AM--GM bound is tight exactly when all $\lambda_i$ are
equal (isotropic block); it loosens as the source power concentrates on fewer
coordinates, since that lowers $\det(I+A)$ at fixed trace.
\end{proof}

\begin{theorem}[Leakage certificate and identification floor]
\label{th:certificate}
Grant Assumptions~\ref{as:range}--\ref{as:indep}, let the private template
$B=f(z_P)$ take $K$ values under a \emph{uniform} prior, and let the adversary
observe $Y=T(z+N)$ with $T$ invertible (the worst case). Define the joint
leakage
\begin{equation}
I_{\mathrm{leak}}(\sigma)
=\tfrac12\!\!\sum_{h\in P}\! d_h\log\!\left(1+\frac{g_h}{d_h\,\sigma_h^{2}}\right),
\label{eq:leak}
\end{equation}
where the bit budget
\begin{equation}
I_{\max}(\tau)=\log_2 K-1-\tau\log_2(K-1)\quad[\text{bit}]
\label{eq:imax}
\end{equation}
is the largest leakage consistent with a floor $\tau$. Here $I_{\max}$ is in
bits ($\log_2$); the allocation program below accumulates mutual information in
nats ($\ln$), so the budget must be converted by $B=I_{\max}\ln2$ before the
KKT solve (Remark~\ref{rm:units}). If
$I_{\mathrm{leak}}(\sigma)\le I_{\max}(\tau)$, then \emph{every} estimator
$\widehat B(Y)$ has identification error $\Pr[\widehat B\neq B]\ge\tau$. This
is an operational attack-error lower bound; it is reported in parallel with the
$(\varepsilon,\delta)$ account of Section~\ref{sec:rdpgdp}, and the two do not
imply one another (Remark~\ref{rm:twosystems}). The feasible floors satisfy
$\tau<\tau_{\max}(K)=(\log_2K-1)/\log_2(K-1)$.
The uniform prior is the strongest case: for a non-uniform prior the entropy
term is $H(B)\le\log_2K$, and the floor weakens by $\log_2K-H(B)$. In
deployment $H(B)$ is estimated from the enrollment prior; the uniform
assumption here treats the gallery as balanced and gives the worst-case prior,
so a measured prior should replace $\log_2K$ in the Fano term.
\end{theorem}

\begin{proof}
The Markov chain $B\to z_P\to Y=T(z+N)\to z+N$ (the last step is the
deterministic $T^{-1}$) gives, by data processing,
$I(B;Y)\le I(z_P;z+N)$. Under the block-wise second-order
structure of Assumption~\ref{as:indep} and Lemma~\ref{lm:maxent},
\begin{equation}
I(z_P;z+N)\le\tfrac12\!\!\sum_{h\in P}\! d_h\log\!\left(1+\frac{g_h}{d_h\,\sigma_h^{2}}\right)
= I_{\mathrm{leak}}(\sigma),
\label{eq:cert-bound}
\end{equation}
the right-hand side of \eqref{eq:leak}. Fano's inequality for a $K$-ary
$B$~\cite{fano} gives
\begin{equation}
\Pr[\widehat B\neq B]\ge\frac{H(B\mid Y)-1}{\log_2(K-1)}
\ge\frac{\log_2 K-I(B;Y)-1}{\log_2(K-1)},
\end{equation}
so $I(B;Y)\le I_{\max}(\tau)$ implies the claimed floor. Sending
$I_{\mathrm{leak}}\downarrow0$ in \eqref{eq:imax} yields the feasibility
condition
\begin{equation}
\tau<\tau_{\max}(K)=\frac{\log_2K-1}{\log_2(K-1)}.
\label{eq:taumax}
\end{equation}
$\tau\ge\tau_{\max}$ would require a negative information budget, which is
infeasible.
\end{proof}

\begin{corollary}[Numbers]
\label{co:numbers}
$\tau_{\max}(10^{4})=0.9248$, $\tau_{\max}(10^{6})=0.9498$, and
$\tau_{\max}(K)\uparrow1$ only logarithmically: an arbitrarily strong
\emph{certified} identification guarantee is unattainable for any finite
template bank, because $H(B)=\log_2K$ is itself the currency. In our
calibration, a floor $\tau=0.90$ with $K=10^{4}$ leaves
$I_{\max}=0.329$\,bit, the entire budget for all private blocks jointly.
Figure~\ref{fig:imax} reports the sensitivity of this budget to the gallery
size: both $\tau_{\max}(K)$ and $I_{\max}$ grow only logarithmically in $K$.
\end{corollary}

\begin{remark}[Nats--bits conversion]
\label{rm:units}
The certificate budget $I_{\max}$ in \eqref{eq:imax} is stated in bits
($\log_2$), whereas the allocation program of Section~\ref{sec:allocation}
accumulates the AWGN mutual information in nats ($\ln$). The two are related
by
\begin{equation}
B = I_{\max}(\tau)\,\ln 2,
\label{eq:natsbits}
\end{equation}
so the $0.329$-bit budget of Corollary~\ref{co:numbers} corresponds to
$B=0.329\ln 2\approx0.228$ nats; the implementation in the accompanying code
applies this conversion before solving the KKT system.
\end{remark}

\begin{remark}[Certificate versus the $(\varepsilon,\delta)$ budget]
\label{rm:twosystems}
Two accounting systems run in parallel here, and both should be reported. The
RDP/GDP/$\ell_1$ machinery of Section~\ref{sec:rdpgdp} turns the accounting rule and noise scale into an
$(\varepsilon,\delta)$ guarantee relative to adjacent records. The Fano certificate of Theorem~\ref{th:certificate} takes a
different route: it needs no $(\varepsilon,\delta)$ budget at all. From the second-order quantities $(g_h,\sigma_h)$
alone it converts the noise level into an operational lower bound on an identification
attack's error. Neither substitutes for the other. A release can meet an $(\varepsilon,\delta)$-DP budget
with a loose $\varepsilon$ yet still leave a measurable identifying trace in the task-relevant
signal. The reverse also occurs: a coarse release that simply drops the identity block
can hold a Fano floor without satisfying DP under the chosen adjacency, since adjacent
records may still differ in some other, noisy block. Tables~\ref{tab:synthetic}--\ref{tab:audit} therefore report
the DP accounting columns and the certificate floor side by side. A like-for-like $(\varepsilon,\delta)$
audit of~\cite{arXiv2510} cannot be made from its reported parameters. It gives a nominal Laplace
scale $\Delta f/\varepsilon$, but the pattern actually released is discriminator-shaped and only
approximately Laplace, with no proven DP guarantee. What follows is therefore a
Gaussian-proxy what-if: we re-price the operating point as if the release were an exact
additive Gaussian mechanism. That is a statement about the accounting law, not an audit
of the learned pattern's true distribution.
\end{remark}

\begin{remark}[Scope of the certificate]
\label{rm:scope}
Theorem~\ref{th:certificate} certifies the identification family, not
arbitrary inference: it constrains every estimator of $B$, but says nothing
about an attribute $A$ that is not a function of $z_P$. It is a task-scoped lower bound derived from second-order statistics and is
logically incomparable with $(\varepsilon,\delta)$-DP under the stated
adjacency (Remark~\ref{rm:twosystems}); it is also a certificate about an
adversary, not a prediction of any measured attack accuracy. Both systems must
be reported, and neither may be presented as the other.
\end{remark}

\begin{remark}[Residual correlation across blocks]
\label{rm:corr}
Hadamard/Fischer's inequality for a PSD block matrix gives
$\det(I+\Sigma_{\mathrm{joint}})\le\prod_h\det(I+\Sigma_{hh})$. At fixed
marginal block variances, off-diagonal correlation decreases the joint
determinant, hence decreases the true joint mutual information; the
block-wise marginal bound in \eqref{eq:cert-bound} is therefore a
\emph{conservative upper bound}, not a lower one. Residual correlation only
widens the Hadamard gap and loosens the bound, but it does not invalidate the
upper bound. A joint covariance estimate would be needed to tighten it; charging
residual correlation to the block source powers $g_h$ is the conservative way to
retain soundness without that estimate.
\end{remark}

% =====================================================================
\section{THE ALLOCATION LAW INDUCED BY THE ACCOUNTING RULE}
\label{sec:allocation}
% =====================================================================

Write $x_h=\sigma_h^{2}$ and $a_h=w_h\kappa_h>0$, where $w_h$ is the
calibrated sensitivity of the task loss to the block and $\kappa_h$ maps
latent power to distortion. The transmitter solves
\begin{equation}\begin{split}
\min_{x_h\ge0}\ &\sum_{h=1}^{H} a_h x_h\\
\text{s.t.}\ &\sum_{h\in P}\varphi_h(x_h)\le B,\quad
\varphi_h(x)=\tfrac{d_h}{2}\log\!\left(1+\frac{g_h}{d_h\,x}\right),
\end{split}\label{eq:primal}
\end{equation}
with $B$ the allowed leakage in nats. Noise is placed only on $P$: a block in
$Q$ carries no obligation and no benefit, so its optimum is $x_h=0$.

\begin{theorem}[Convexity, uniqueness, and closed form]
\label{th:allocation}
The problem \eqref{eq:primal} is a convex program with a unique optimum when
$B\in(0,\infty)$. Let $\nu\ge0$ denote its dual multiplier. The KKT solution is
\begin{equation}
x_h^{\star}(\nu)=\frac{g_h}{2d_h}\left(
\sqrt{1+\frac{2\nu\,d_h^{2}}{a_h g_h}}-1\right),\qquad h\in P,
\label{eq:kkt}
\end{equation}
pinned by $\sum_{h\in P}\varphi_h(x_h^\star(\nu))=B$. The map
$\nu\mapsto\sum_{h\in P}\varphi_h(x_h^\star(\nu))$ is strictly decreasing and
continuous from $+\infty$ ($\nu\downarrow0$) to $0$ ($\nu\uparrow\infty$), so
the bisection in the accompanying code is well posed and converges
geometrically in $\log\nu$. The dual multiplier $\nu$ has a direct physical
reading: it is the shadow price of the leakage budget, i.e.\ the marginal
utility distortion charged for releasing one additional unit of private
information. A tight budget drives $\nu$ up: every private block is asked for
more noise, and the marginal cost of further leakage rises.
\end{theorem}

\begin{proof}
$\varphi_h$ is convex on $x>0$:
$\varphi_h'(x)=-\frac{d_hg_h}{2x(d_hx+g_h)}$ and
$\varphi_h''(x)=\frac{d_hg_h(2d_hx+g_h)}{2x^2(d_hx+g_h)^2}>0$; the objective is
linear, so the feasible set is convex and the program is convex. At an optimum
the constraint is active because $B>0$ forces some $x_h<\infty$, and on the
active set every $x_h>0$ for $h\in P$; strict convexity of $\varphi_h$ then
gives a unique minimiser. Stationarity reads
$a_h+\nu\,\varphi_h'(x_h)=0$, i.e.
$a_h=\nu d_hg_h/(2x_h(d_hx_h+g_h))$, which rearranges to
$x_h^2+(g_h/d_h)x_h-\nu g_h/(2a_h)=0$; its positive root is \eqref{eq:kkt}.
Monotonicity: $\partial x_h^\star/\partial\nu>0$ for every $h$, and
$\varphi_h$ is decreasing in $x$, so each summand decreases strictly; the
limits follow from $x_h^\star\downarrow0$ and $x_h^\star\uparrow\infty$. A
continuous strictly monotone scalar function is inverted by bisection, and
each halving of the $\log\nu$ interval halves the bracketing width, whence
geometric convergence. Each bisection step evaluates the KKT system over the
$|P|$ private blocks, so the allocation costs $O(|P|\,L)$ for $L$ bisection
steps; the later RDP conversion adds $O(|P|\,A)$ over $A$ candidate orders.
Both are offline computations over a small block count and do not constrain the
release rate.
\end{proof}

\begin{corollary}[Two asymptotics]
\label{co:asymptotics}
From \eqref{eq:kkt},
\begin{align}
x_h^\star&\simeq\frac{\nu d_h}{2a_h}, & x_h&\ll g_h/d_h,\label{eq:asym1}\\
x_h^\star&\simeq\sqrt{\frac{\nu g_h}{2a_h}}, & x_h&\gg g_h/d_h.\label{eq:asym2}
\end{align}
The first regime is task-driven and independent of $g_h$: the noise variance
stays below the block source power, and allocation follows the task weight
$a_h$. This corresponds to private blocks that remain close to the task
subspace, where Bob still needs the block and the allowed noise is held to the
task's tolerance. The second is geometry-driven and grows only as $\sqrt{g_h}$:
the noise variance exceeds the source power, and blocks carrying large private
power but low task cost absorb most of the noise. A component that is cheap for
the task but expensive to protect ($a_h$ small, $g_h$ large) absorbs most of
the noise precisely in this second regime.
\end{corollary}

\begin{theorem}[The accounting rule dictates the shape]
\label{th:shape}
Replace the leakage constraint of \eqref{eq:primal} by each of the standard
accounting laws at a common budget $\varepsilon$, and define
\begin{equation}
c=\sqrt{2\log(1.25/\delta)}.
\label{eq:tailconst}
\end{equation}
The stationary profiles are:
\begin{enumerate}
\item[(a)] \emph{amplitude-only $\ell_1$ composition}
$\sum_h c\Delta f_h/\sigma_h\le\varepsilon$:
\begin{equation}
x_h^{\star}\propto\Big(\frac{\Delta f_h}{a_h}\Big)^{2/3},
\label{eq:l1shape}
\end{equation}
a geometric profile, strongly unequal and dominated by components with large
sensitivity and small task weight.
\item[(b)] \emph{quadrature composition} (RDP and GDP alike,
Section~\ref{sec:rdpgdp}): the constraint is $\sum_h\chi_h/x_h\le B'$ with
$\chi_h=\Delta f_h^2$ (GDP) or $\alpha\Delta f_h^2/2$ (RDP at order $\alpha$),
so stationarity $a_h=\nu\chi_h/x_h^2$ gives
\begin{equation}
x_h^{\star}\propto\frac{\Delta f_h}{\sqrt{a_h}},
\label{eq:quadshape}
\end{equation}
i.e.\ exponent $1/2$ in the task weight. This is the same profile as the rigorous
tail bound at large $\mu$ (part (c)), and \emph{not} the logarithmic profile
\eqref{eq:kkt}. The two coincide only in the geometry-driven limit $x\gg g$ of
the log constraint, where substituting $g_h\mapsto\Delta f_h^2$ makes the
quadratic terms agree.
\item[(c)] \emph{rigorous tail bound} (Equation~\eqref{eq:l1exact} below):
\begin{equation}
2a_hx_h^{2}=\nu\big(\Delta f_h^{2}+c\Delta f_h\sqrt{x_h}\big),
\label{eq:l1exactshape}
\end{equation}
so $x_h\propto\Delta f_h/\sqrt{a_h}$ when $\mu_h\gg c$ (the quadratic term
dominates) and \eqref{eq:l1shape} when $\mu_h\ll c$ (the linear term dominates),
with $\mu_h=\Delta f_h/\sigma_h$.
\end{enumerate}
The three stationary profiles are plotted side by side in
Figure~\ref{fig:budget_split}.
\end{theorem}

\begin{proof}
(a) The constraint is $\sum_hb_hx_h^{-1/2}\le\varepsilon$ with
$b_h=c\Delta f_h$; it is convex since $\psi_h(x)=b_hx^{-1/2}$ has
$\psi_h''>0$. Stationarity $a_h=\nu b_hx_h^{-3/2}/2$ gives
$x_h=(2a_h/(\nu b_h))^{-2/3}\propto(b_h/a_h)^{2/3}$, which is
\eqref{eq:l1shape}; $\nu$ is fixed by the budget.
(b) GDP accumulates in quadrature,
$\sum_h(\Delta f_h/\sigma_h)^{2}\le c_{\mathrm{GDP}}(\varepsilon)$ by
\eqref{eq:gdpsum} inverted in $\varepsilon$; RDP accumulates
$\sum_h\alpha\Delta f_h^2/(2\sigma_h^2)\le\varepsilon$ at the selected order
$\alpha$ by \eqref{eq:rdpsum}. Both are of the form $\sum_h\chi_h/x_h\le B'$
with $\chi_h\propto\Delta f_h^2$. Stationarity of
$L=\sum_ha_hx_h+\nu(\sum_h\chi_h/x_h-B')$ gives
$a_h-\nu\chi_h/x_h^2=0$, hence $x_h=\sqrt{\nu\chi_h/a_h}\propto\Delta f_h/\sqrt{a_h}$,
which is \eqref{eq:quadshape}. This is \emph{not} the KKT system of the log-leakage
program \eqref{eq:primal}: that system yields exponent $1$ only in its small-$x$
(task-driven) regime and exponent $1/2$ in the large-$x$ (geometry-driven)
regime, and recovers \eqref{eq:quadshape} solely when $g_h$ is replaced by
$\Delta f_h^2$ in that limit.
(c) The constraint \eqref{eq:l1exact} has
$f_h(x)=\Delta f_h^2/(2x)+c\Delta f_hx^{-1/2}$; stationarity
$a_h=\nu f_h'(x_h)$ reads
$2a_hx_h^2=\nu\Delta f_h(\Delta f_h+c\sqrt{x_h})$, which is
\eqref{eq:l1exactshape}. Dividing by $c\Delta f_h\sqrt{x_h}$ identifies the
governing ratio $\mu_h/c=\Delta f_h/(c\sqrt{x_h})$: for $\mu_h\gg c$ the term
$\Delta f_h^2$ dominates and $x_h^2\approx\nu\Delta f_h^2/(2a_h)$, giving
$x_h\propto\Delta f_h/\sqrt{a_h}$; for $\mu_h\ll c$ the term
$c\Delta f_h\sqrt{x_h}$ dominates and \eqref{eq:l1shape} is recovered.
\end{proof}

\begin{remark}[Two exponents, four laws, one message]
\label{rm:exponents}
Across the allocation laws the exponent of the task weight $a_h$ takes only two
distinct values. The logarithmically-metered certificate program
\eqref{eq:primal} has exponent $1$ in its task-driven (small-$x$) regime and
$1/2$ in its geometry-driven (large-$x$) regime; the quadrature laws
\eqref{eq:quadshape} (RDP/GDP) and the rigorous tail bound at large $\mu$
(Theorem~\ref{th:shape}(c)) both use exponent $1/2$; the amplitude-only
$\ell_1$ law uses $2/3$, and the rigorous tail at small $\mu$ reverts to $2/3$.
All laws agree that noise belongs where $a_h$ is small; they differ in how
strongly they push. The static-index policy has no exponent at all, because it
cannot use the $a_h$ scale.
\end{remark}

\begin{table}[h]
\centering\footnotesize
\setlength{\tabcolsep}{3.5pt}
\caption{Applicability of the three accounting laws to a Gaussian additive
release.}
\label{tab:laws}
\begin{tabular}{p{2.5cm}p{2.6cm}p{2.6cm}}
\toprule
law & composition form & when it applies \\
\midrule
Amplitude-only $\ell_1$ & $\sum_h c\Delta f_h/\sigma_h$ & Valid only when
$\mu_h\lesssim c$; under-reports $\varepsilon$ at large $\mu$. \\
Quadrature (RDP/GDP) & $\sum_h(\Delta f_h/\sigma_h)^{2}$ & Correct account
of the Gaussian additive mechanism; the recommended standard. \\
Rigorous linear tail & $\sum_h(\mu_h^2/2+c\mu_h)$ & Valid at every $\mu_h$,
but conservative; lies above the quadrature account. \\
\bottomrule
\end{tabular}
\end{table}

\begin{corollary}[Static index selection is a corner of this program]
\label{co:static}
The reference strategy, which protects a fixed set $S$ of $k$ codes with a common
noise level and leaves the rest untouched, is obtained by imposing on the
$\ell_1$ program the extra constraint
\begin{equation}
x_h=
\begin{cases}
\sigma_{\max}^{2}, & h\in S,\\
\sigma_{\min}^{2}\approx 0, & h\notin S,
\end{cases}
\qquad |S|=k.
\end{equation}
Because the untouched private blocks are released at $\sigma_{\min}^2\approx0$, the
$\ell_1$ leakage $\sum_h c\Delta f_h/\sigma_h$ diverges: this point is \emph{outside}
the feasible set of the budgeted program, not a corner of it. It is therefore
not a comparable feasible policy: it exposes the clean private blocks at zero
cost in distortion while violating the privacy constraint outright, which is
why its certified floor collapses to zero (Section~\ref{sec:experiments}). The
static rule cannot outperform \eqref{eq:l1shape} on any common feasible
budget, and it is strictly worse whenever the $\ell_1$ optimum is supported on
more than $k$ components with unequal amplitudes, the generic case since
$\Delta f_h$ and $a_h$ are heterogeneous. Dynamic allocation replaces the
combinatorial choice of $S$ by a global multiplier.
\end{corollary}

% =====================================================================
\section{RDP, GAUSSIAN-DP, AND THEIR COMPOSITION}
\label{sec:rdpgdp}
% =====================================================================

\emph{Gaussian mechanism, one block.} Let $q(z)=z$ have sensitivity
$\Delta_h=\Delta f_h$ on the block. For neighbouring $z,z'$ with
$\delta_h=z'-z$, $\|\delta_h\|_2\le\Delta_h$, the laws of $Y_h$ are
$\mathcal{N}(z_h,\sigma_h^2I)$ and $\mathcal{N}(z_h',\sigma_h^2I)$. The
Rényi divergence of order $\alpha>1$ between two such Gaussians is
\begin{equation}
\begin{aligned}
D_\alpha\!\left(\mathcal{N}(z_h,\sigma_h^2I)\,\|\,
\mathcal{N}(z_h',\sigma_h^2 I)\right)
&=\frac{\alpha\|\delta_h\|^2}{2\sigma_h^2}
\le\frac{\alpha\Delta f_h^{2}}{2\sigma_h^{2}}\\
&=:\varepsilon_h^{\RDP}(\alpha),
\end{aligned}
\label{eq:rdpblock}
\end{equation}
obtained by evaluating the $\alpha$-th moment of the density ratio under the
shifted law.

\emph{Composition (RDP).} Rényi divergence adds over independent mechanism
outputs~\cite{mironov2017rdp}; the block releases are independent by
\eqref{eq:mechanism}, so
\begin{equation}
\varepsilon^{\RDP}(\alpha)=\sum_{h\in P}\varepsilon_h^{\RDP}(\alpha)
=\frac{\alpha}{2}\sum_{h\in P}\frac{\Delta f_h^{2}}{\sigma_h^{2}}.
\label{eq:rdpsum}
\end{equation}
Equation~\eqref{eq:rdpsum} is the quadrature, not the $\ell_1$, accumulation,
and grouping of codes does not change it. The composition order matters: the
per-block RDP divergences are summed first at each order $\alpha$, and the
single summed RDP value is then converted to $(\varepsilon,\delta)$; one does
not convert each block to an $(\varepsilon,\delta)$ pair and add the resulting
$\varepsilon$ values, which would reintroduce a loose linear tail account. The
GDP account below follows the same order: the per-block $\mu_h$ are combined in
quadrature first, and the single total $\mu_{\mathrm{tot}}$ is then inverted
through the Gaussian privacy-loss law.

\emph{Conversion to $(\varepsilon,\delta)$.} An
$(\alpha,\varepsilon^{\RDP})$-RDP mechanism is
$(\varepsilon^{\RDP}+\log(1/\delta)/(\alpha-1),\delta)$-DP for every
$\alpha>1$~\cite{mironov2017rdp}; optimising over $\alpha$,
\begin{equation}
\varepsilon(\delta)=\min_{\alpha>1}\left[
\frac{\alpha}{2}\sum_{h\in P}\frac{\Delta f_h^{2}}{\sigma_h^{2}}
+\frac{\log(1/\delta)}{\alpha-1}\right].
\label{eq:rdpconv}
\end{equation}
The bracket is convex in $\alpha$, so the minimiser is unique; the
implementation in the accompanying code searches over an order grid. The
conversion is conservative: tighter conversions are available in the
hypothesis-testing privacy literature, and \eqref{eq:rdpconv} is kept so that
all reported $\varepsilon$ are upper bounds.

\emph{Privacy loss random variable and Gaussian-DP.} For one block the privacy
loss of an outcome $y$ is
\begin{equation}
L_h(y)=\log\frac{p_{z_h}(y)}{p_{z_h'}(y)}
=\frac{\|\delta_h\|^{2}-2\langle y-z_h,\delta_h\rangle}{2\sigma_h^{2}},
\label{eq:loss}
\end{equation}
using $p_{z'}(y)=p_z(y-\delta_h)$. Under $y\sim p_{z_h}$,
$u=y-z_h$ has mean zero and covariance $\sigma_h^2I$, so
\begin{equation}
L_h\sim\mathcal{N}\!\left(\frac{\mu_h^{2}}{2},\,\mu_h^{2}\right),
\qquad \mu_h:=\frac{\Delta f_h}{\sigma_h}.
\label{eq:gdpblock}
\end{equation}
The summands are independent, so two Gaussians compose into a Gaussian:
\begin{equation}\begin{split}
L&=\sum_{h\in P}L_h
\sim\mathcal{N}\!\left(\frac{\mu_{\mathrm{tot}}^{2}}{2},
\mu_{\mathrm{tot}}^{2}\right),\\
\mu_{\mathrm{tot}}&=\sqrt{\sum_{h\in P}\mu_h^{2}}
=\sqrt{\sum_{h\in P}\frac{\Delta f_h^{2}}{\sigma_h^{2}}}.
\end{split}\label{eq:gdpsum}
\end{equation}
This is Gaussian differential privacy~\cite{dong2022gdp}. For the Gaussian
mechanism the privacy-loss random variable is itself Gaussian, so this GDP
characterization is \emph{exact} for the additive Gaussian release rather than
an approximation; the quadrature composition below inherits that exactness.
The conservative
one-sided form used in the code and in all tables is
\begin{equation}
\delta(\varepsilon)\le\Phi\!\left(-\frac{\varepsilon}{\mu_{\mathrm{tot}}}
+\frac{\mu_{\mathrm{tot}}}{2}\right),
\label{eq:gdpdelta}
\end{equation}
whose inverse is
$\varepsilon(\delta)=\mu_{\mathrm{tot}}(\mu_{\mathrm{tot}}/2-\Phi^{-1}(\delta))$.
The exact two-sided expression subtracts the second tail term; it differs
from \eqref{eq:gdpdelta} only for $\varepsilon$ of order
$\mu_{\mathrm{tot}}^2/2$, where both forms are already vacuous, so
\eqref{eq:gdpdelta} is the standard choice for tables.

\emph{Linear composition (basic accounting).} The classical Gaussian-mechanism
tail bound, composed linearly~\cite{dwork2014}, keeps the quadratic privacy
term:
\begin{equation}
\varepsilon_{\ell_1}(\delta)=\sum_{h\in P}\Big(\mu_h^{2}/2+c\,\mu_h\Big),
\qquad \mu_h=\frac{\Delta f_h}{\sigma_h},
\label{eq:l1exact}
\end{equation}
with $c$ of \eqref{eq:tailconst}. The amplitude-only truncation
$\sum_hc\Delta f_h/\sigma_h$ drops the $\mu_h^{2}/2$ term; it is the law that
Laplace-style accounting effectively assumes. The dropped term is not cosmetic.
At the low-noise (large-$\mu$) operating point of the reference
scheme~\cite{arXiv2510},
$\Delta f_h\approx66$--$100$ (representative $\approx91$), $\sigma_h\approx3.6$, hence $\mu_h\approx24.1$ (RMS) and
$|P|=6$, and $\mu_h\gg c=4.84$;
\eqref{eq:l1exact} is then dominated by $\mu_h^2/2$. On the
synthetic anchor at the release certified at quadrature $\varepsilon=2000$,
$\delta=10^{-5}$, with $|P|=6$ and $\mu_h\approx24.1$, the amplitude-only
truncation reports $\varepsilon\approx692.6$, while
\eqref{eq:l1exact} reports $\approx2440.5$; the amplitude-only figure
\emph{under-reports} the loss and is unsoundly tighter-looking than the Gaussian
tail account. The same audit on
the measured proxy (Table~\ref{tab:audit}) gives $701.1$ and $2448.9$, with the
same ordering; the three abstract figures $693/2441/2000$ are the anchor point,
and $2000$ is the quadrature certificate they are priced against.
Figure~\ref{fig:composite}(b) plots this audit over the whole range. This is
the limitation of a ``the noise resembles Laplace, hence
$\varepsilon=\Delta f/s$'' argument: the ratio form is dimensionally plausible
but, for a Gaussian release, omits the quadratic privacy-loss term. Applying a
Laplace-style $\ell_1$ composition to this
Gaussian release at a low-noise operating point yields an assessment that looks
tighter than the mechanism actually guarantees; the reported budget should
therefore be read against the quadrature or tail account, not the amplitude-only
ratio.

\begin{proposition}[Quadrature composition is the correct Gaussian account]
\label{pr:quadrature}
For the additive Gaussian release of \eqref{eq:mechanism}, the correct
privacy-loss composition is the quadratic (orthogonal) combination
$\sum_h\mu_h^{2}$, not a linear $\ell_1$ sum $\sum_h\mu_h$. Formally, for any
$\{\sigma_h\}$ and any $P$,
$\sum_{h\in P}\Delta f_h^{2}/\sigma_h^{2}
\le\big(\sum_{h\in P}\Delta f_h/\sigma_h\big)^{2}$, with equality iff at most
one summand is non-zero. At the divergence level $\sum\mu_h^2\le(\sum\mu_h)^2$, but after conversion to
$(\varepsilon,\delta)$ the quadratic privacy-loss term $\mu_h^2/2$ makes the
quadrature $\varepsilon$ larger than the amplitude-only figure at low noise
(Table~\ref{tab:audit}): quadrature is the \emph{correct} account, not a
numerically smaller one. It is the accumulation law that matches the actual
additive-noise mechanism of \eqref{eq:mechanism}.
\end{proposition}

\begin{proof}
The inequality is $\|\cdot\|_2^2\le\|\cdot\|_1^2$ on the vector
$(\Delta f_h/\sigma_h)_{h\in P}$, with equality iff the vector is supported on
a single coordinate. The last claim is \eqref{eq:rdpsum}--\eqref{eq:gdpsum}
computed from the likelihood ratio of the actual release.
\end{proof}

% =====================================================================
%  figures_pgfplots.tex  --  camera-ready vector figures for the
%  privacy-accounting section.  Every plot reads a .dat file written by
%  `python -m sdpa.figures`, so text, ticks and fonts are LaTeX-native.
%
%  Set the data directory before \input, e.g.
%      \def\figdatadir{../figures/dat/}
%  The default below assumes you compile from inside figures/.
% =====================================================================
\providecommand{\figdatadir}{dat/}
\pgfplotsset{compat=1.18}

% ---------------------------------------------------------------------
% Fig. 1 -- the accounting law dictates the noise profile
% ---------------------------------------------------------------------
\begin{figure}[t]
\centering
\begin{tikzpicture}
  \pgfplotstableread{\figdatadir budget_split.dat}\bs
  \begin{axis}[
      width=\columnwidth, height=0.70\columnwidth,
      xlabel={semantic component $h$},
      ylabel={share of noise power},
      xmin=-0.8, xmax=14.4, ymin=0, ymax=0.30,
      xtick={0,2,...,14},
      tick label style={font=\footnotesize},
      label style={font=\footnotesize},
      legend style={font=\scriptsize, draw=none, fill=none,
                    at={(0.5,-0.26)}, anchor=north, legend columns=3},
      ymajorgrids, grid style={densely dotted, black!25},
      extra x ticks={0,1,3,4,5,6},
      extra x tick labels={},
      extra tick style={grid=major, major grid style={black!15}},
    ]
    \addplot[draw=none, fill=gray!10] coordinates {
      (1.5,0)(2.5,0)(2.5,0.30)(1.5,0.30)(1.5,0)
      (6.5,0)(7.5,0)(7.5,0.30)(6.5,0.30)(6.5,0)
      (7.5,0)(8.5,0)(8.5,0.30)(7.5,0.30)(7.5,0)
      (8.5,0)(9.5,0)(9.5,0.30)(8.5,0.30)(8.5,0)
      (9.5,0)(10.5,0)(10.5,0.30)(9.5,0.30)(9.5,0)
      (10.5,0)(11.5,0)(11.5,0.30)(10.5,0.30)(10.5,0)
      (11.5,0)(12.5,0)(12.5,0.30)(11.5,0.30)(11.5,0)
      (12.5,0)(13.5,0)(13.5,0.30)(12.5,0.30)(12.5,0)
      (13.5,0)(14.5,0)(14.5,0.30)(13.5,0.30)(13.5,0)};
    \addplot+[ybar, bar width=4pt, fill=blue!55!black, draw=none]
      table[x expr=\thisrow{h}-0.27, y=gdp] {\bs};
    \addplot+[ybar, bar width=4pt, fill=blue!30, draw=none]
      table[x expr=\thisrow{h}-0.09, y=rdp] {\bs};
    \addplot+[ybar, bar width=4pt, fill=orange!70!black, draw=none]
      table[x expr=\thisrow{h}+0.09, y=basic] {\bs};
    \addplot+[ybar, bar width=4pt, fill=black!40, draw=none]
      table[x expr=\thisrow{h}+0.27, y=static] {\bs};
    \addplot[black, mark=*, mark size=1pt, densely dotted, mark options={solid}]
      table[x=h, y=wkappa_norm] {\bs};
    \legend{GDP, RDP, $\ell_1$ (ampl.), static ($k{=}4$), $w_h\kappa_h$}
  \end{axis}
\end{tikzpicture}
\caption{Noise-power share of the four accounting rules on the task-decoupled
calibration ($\varepsilon{=}100$, $\delta{=}10^{-5}$; this panel illustrates
the shape of the allocation, while the audit in Fig.~\ref{fig:composite}(b)
uses the $\varepsilon{=}2000$ operating point). Only the declared private
components $h\in P=\{0,1,3,4,5,6\}$ may carry noise
(Theorem~\ref{th:allocation}); the gray-shaded blocks are the task complement
$Q$. The amplitude-only $\ell_1$ rule concentrates noise on a few high-sensitivity
components, the quadrature GDP/RDP rules spread it more evenly, and the static
rule adds noise to only the four selected blocks and leaves the rest untouched
(Corollary~\ref{co:static}).}
\label{fig:budget_split}
\end{figure}

% ---------------------------------------------------------------------
% Fig. 2 -- the price of a certificate
% ---------------------------------------------------------------------
\begin{figure}[t]
\centering
\begin{tikzpicture}
  \pgfplotstableread{\figdatadir certificate_price.dat}\cp
  \begin{axis}[
      name=left, width=0.84\columnwidth, height=0.50\columnwidth,
      xlabel={certified error floor $\tau$},
      ylabel={Bob LPIPS},
      xmin=0.29, xmax=0.96, ymin=0.10, ymax=1.0,
      tick label style={font=\footnotesize},
      label style={font=\footnotesize},
      legend style={font=\scriptsize, draw=none, fill=none,
                    at={(0.5,-0.24)}, anchor=north, legend columns=2,
                    nodes={font=\scriptsize, inner sep=1pt}},
      ymajorgrids, grid style={densely dotted, black!25},
    ]
    \addplot[red!70!black, mark=square*, mark size=1.2pt]
      table[x=tau, y=lpips_aligned] {\cp};
    \addplot[blue!70!black, mark=*, mark size=1.2pt]
      table[x=tau, y=lpips_decoupled] {\cp};
    \addplot[black, densely dotted, domain=0.29:0.96] {0.112};
    \legend{{Bob, aligned $w_h$}, {Bob, decoupled $w_h$},
            {no-protection $0.112$}}
  \end{axis}
  \begin{axis}[
      name=right, at={(left.south west)}, anchor=south west,
      width=0.84\columnwidth, height=0.50\columnwidth,
      axis y line*=right, axis x line=none, xtick=\empty,
      ylabel={identification error $P_e$},
      ymin=0.25, ymax=1.15,
      tick label style={font=\footnotesize},
      label style={font=\footnotesize},
    ]
    \addplot[blue!70!black, densely dotted, very thick]
      table[x=tau, y=certified_floor] {\cp};
    \addplot[blue!70!black, only marks, mark=triangle*, mark size=1.4pt,
             densely dashed]
      table[x=tau, y=p_e_surrogate] {\cp};
  \end{axis}
\end{tikzpicture}
\caption{Certificate price on the synthetic anchor, two overlap regimes. The
two Bob LPIPS curves are read on the \emph{left} axis; the dotted
certified-floor curve and the triangle surrogate $P_e$ are read on the
\emph{right} axis. The LPIPS values are surrogate predictions from the default
distortion model, not measured reconstruction distortion. When the private
semantics \emph{are} the task (aligned, $w_h$ large, task overlap $\approx1$)
the certified floor drives Bob's LPIPS to $0.63$--$0.87$, leaving little usable
signal for the legitimate receiver; when the private subspace is task-null
(decoupled) the same floor costs almost nothing until
$\tau$ approaches the ceiling $\tau_{\max}(10^{4}){=}0.925$. The dotted line
is the certified floor itself, which the surrogate $P_e$ must dominate
(Theorem~\ref{th:certificate}).}
\label{fig:certificate}
\end{figure}

% ---------------------------------------------------------------------
% Fig. 4 -- leakage budget I_max(tau, K) sensitivity to the gallery size
% ---------------------------------------------------------------------
\begin{figure}[t]
\centering
\begin{tikzpicture}
  \pgfplotstableread{\figdatadir imax_gallery.dat}\ig
  \begin{axis}[
      width=\columnwidth, height=0.62\columnwidth,
      xlabel={certified error floor $\tau$},
      ylabel={$I_{\max}$ (bit)},
      xmin=0.30, xmax=0.92, ymin=0,
      tick label style={font=\footnotesize},
      label style={font=\footnotesize},
      legend style={font=\scriptsize, draw=none, fill=none,
                    at={(0.98,0.98)}, anchor=north east},
      ymajorgrids, grid style={densely dotted, black!25},
    ]
    \addplot[blue!55!black, dashed] table[x=tau, y=k1e3] {\ig};
    \addplot[blue!70!black, mark=*, mark size=1.2pt] table[x=tau, y=k1e4] {\ig};
    \addplot[orange!80!black, dotted, very thick] table[x=tau, y=k1e6] {\ig};
    \legend{$K=10^{3}$, $K=10^{4}$, $K=10^{6}$}
  \end{axis}
\end{tikzpicture}
\caption{Leakage budget $I_{\max}(\tau,K)$ of \eqref{eq:imax} versus the
certified floor, for three identity-gallery sizes. Raising $K$ pushes the
feasibility ceiling $\tau_{\max}(K)$ and the budget upward only
logarithmically: adjacent gallery sizes differ by about $0.025$ in
$\tau_{\max}$ and the full span over three orders of magnitude ($10^{3}$ to
$10^{6}$) is about $0.050$. Curves are truncated at their own feasibility
ceiling $\tau_{\max}(K)$; beyond it no non-negative budget exists.}
\label{fig:imax}
\end{figure}

% ---------------------------------------------------------------------
% Fig. 3 -- composite, one column per story
% ---------------------------------------------------------------------
\begin{figure*}[t]
\centering
\begin{tikzpicture}
  \pgfplotstableread{\figdatadir waterfill_sweep.dat}\wf
  \begin{axis}[
      width=0.315\textwidth, height=0.27\textwidth,
      xlabel={dual variable $\nu$}, ylabel={$\sigma_h^{2}/g_h$},
      xmode=log, ymode=log, ymin=1e-4, ymax=1e8,
      tick label style={font=\footnotesize},
      label style={font=\footnotesize},
      title={(a) water-filling}, title style={font=\footnotesize},
      legend style={font=\scriptsize, draw=none, fill=none,
                    at={(0.98,0.03)}, anchor=south east},
    ]
    \addplot[blue!55!black] table[x=mu, y=w_1] {\wf};
    \addplot[blue!55!black, densely dashed] table[x=mu, y=w_8] {\wf};
    \addplot[orange!80!black, densely dotted, very thick]
      table[x=mu, y=w_125] {\wf};
    \legend{$w_h$, $w_h/8$, $w_h/125$}
  \end{axis}
\end{tikzpicture}
\hfill
\begin{tikzpicture}
  \pgfplotstableread{\figdatadir accounting_audit.dat}\ac
  \begin{axis}[
      width=0.315\textwidth, height=0.27\textwidth,
      xlabel={noise-to-signal ratio (dB)}, ylabel={certified $\epsilon$},
      ymode=log,
      tick label style={font=\footnotesize},
      label style={font=\footnotesize},
      title={(b) audited $\epsilon$ at one release},
      title style={font=\footnotesize},
      legend style={font=\scriptsize, draw=none, fill=none,
                    at={(0.98,0.03)}, anchor=south east},
    ]
    \addplot[draw=none, fill=red!7] coordinates {
      (9.7,1)(24,1)(24,5000)(9.7,5000)(9.7,1)};
    \node[font=\scriptsize, red!60!black, anchor=north west]
      at (axis cs:10.2,3000) {hazard: $\mu_h\gg c$};
    \addplot[blue!55!black, mark=triangle*, mark size=1.5pt]
      table[x=nsr_db, y=eps_quadrature] {\ac};
    \addplot[green!40!black, mark=square*, mark size=1.2pt]
      table[x=nsr_db, y=eps_l1_exact] {\ac};
    \addplot[orange!80!black, mark=*, mark size=1.2pt, densely dashed]
      table[x=nsr_db, y=eps_l1_stylized] {\ac};
    \legend{quadrature, $\ell_1$ exact, $\ell_1$ ampl.}
  \end{axis}
\end{tikzpicture}
\hfill
\begin{tikzpicture}
  \pgfplotstableread{\figdatadir certificate_gap.dat}\cg
  \begin{axis}[
      width=0.315\textwidth, height=0.27\textwidth,
      xlabel={certified floor $\tau$}, ylabel={$P_e^{\rm sim}-\tau$},
      ymin=-0.05, ymax=0.62,
      tick label style={font=\footnotesize},
      label style={font=\footnotesize},
      title={(c) certificate tightness},
      title style={font=\footnotesize},
    ]
    \addplot[purple!60!black, mark=*, mark size=1.2pt]
      table[x=tau, y=gap] {\cg};
    \addplot[black, densely dotted, domain=0.42:0.93] {0};
  \end{axis}
\end{tikzpicture}
\caption{(a) Water-filling charge: one latent direction, three task weights.
Privacy is cheaper for a task-null direction because $a_h=w_h\kappa_h$
enters the KKT solution \eqref{eq:kkt}. (b) Audit of the accounting laws on
the \emph{same} release at the low-noise (large-$\mu$) operating point certified at
$\varepsilon{=}2000$, $\delta{=}10^{-5}$: quadrature (GDP) sits between the
rigorous tail bound and the amplitude-only $\ell_1$ form, the latter under-
reporting $\varepsilon$ once $\mu_h\gg c$. The red shaded band marks this
large-$\mu$ hazard region, where a Laplace-style $\ell_1$ composition yields an
assessment tighter than the guarantee. (c) The surrogate $P_e$ dominates the
certified floor everywhere: the certificate is conservative, not optimistic.}
\label{fig:composite}
\end{figure*}

% ---------------------------------------------------------------------
% Fig. 4 -- latent-level measurement on the released SemanticStyleGAN
% ---------------------------------------------------------------------
\begin{figure*}[t]
\centering
\begin{tikzpicture}
  \pgfplotstableread{\figdatadir latent_utility_landscape.dat}\lu
  \begin{axis}[
      width=0.315\textwidth, height=0.27\textwidth,
      ybar, bar width=3.2pt,
      ylabel={measured $w_h$},
      ymin=0, ymax=0.40,
      xtick={0,...,14},
      xticklabels={sh0,sh1,bg,skin,eye,brw,mth,nos,ear,hair,nck,clth,eyg,hat,erR},
      x tick label style={rotate=55, anchor=east, font=\tiny},
      tick label style={font=\footnotesize},
      label style={font=\footnotesize},
      title={(a) measured utility landscape}, title style={font=\footnotesize},
    ]
    \addplot[fill=blue!45, draw=blue!55!black] table[x=h, y=w_h] {\lu};
  \end{axis}
\end{tikzpicture}
\hfill
\begin{tikzpicture}
  \pgfplotstableread{\figdatadir latent_sweep_examples.dat}\ls
  \begin{axis}[
      width=0.315\textwidth, height=0.27\textwidth,
      xlabel={injected $\sigma_h^{2}$}, ylabel={Bob LPIPS},
      xmode=log,
      tick label style={font=\footnotesize},
      label style={font=\footnotesize},
      title={(b) per-block noise sweeps}, title style={font=\footnotesize},
      legend style={font=\scriptsize, draw=none, fill=none,
                    at={(0.03,0.97)}, anchor=north west},
    ]
    \addplot[blue!55!black, mark=*, mark size=1.2pt]
      table[x=sigma2, y=shared0] {\ls};
    \addplot[orange!80!black, mark=square*, mark size=1.2pt]
      table[x=sigma2, y=skin] {\ls};
    \addplot[green!40!black, mark=triangle*, mark size=1.5pt]
      table[x=sigma2, y=hair] {\ls};
    \addplot[purple!60!black, mark=diamond*, mark size=1.5pt, densely dashed]
      table[x=sigma2, y=eyes] {\ls};
    \legend{shared\_0, skin, hair, eyes}
  \end{axis}
\end{tikzpicture}
\hfill
\begin{tikzpicture}
  \pgfplotstableread{\figdatadir latent_kappa.dat}\lk
  \begin{axis}[
      width=0.315\textwidth, height=0.27\textwidth,
      xlabel={injected $\sigma_h^{2}$}, ylabel={residual fraction $\kappa_h$},
      xmode=log, ymin=0, ymax=1.05,
      tick label style={font=\footnotesize},
      label style={font=\footnotesize},
      title={(c) Wiener deprotector residual}, title style={font=\footnotesize},
      legend style={font=\scriptsize, draw=none, fill=none,
                    at={(0.97,0.97)}, anchor=north east},
    ]
    \addplot[blue!55!black, mark=*, mark size=1.2pt]
      table[x=sigma2, y=kappa_empirical] {\lk};
    \addplot[red!70!black, densely dashed, very thick]
      table[x=sigma2, y=kappa_analytic] {\lk};
    \legend{empirical (held-out), analytic}
  \end{axis}
\end{tikzpicture}
\caption{Latent-level measurement on the released SemanticStyleGAN generator.
(a) Measured per-block utility weights $w_h$: the perceptual cost of latent
noise concentrates on the shared codes, hair, skin and background, while the
facial-feature blocks (eyes, nose, ear\_r) are nearly free to protect. Block
order: shared\_0, shared\_1, background, skin, eyes, brows, mouth, nose, ears,
hair, neck, cloth, eye\_g, hat, ear\_r. (b) Example decode-and-LPIPS sweeps
with the quadratic fit $\mathrm{LPIPS}\approx\mathrm{lpips}_0+w_h\sigma_h^2$;
the eyes curve is almost flat---the generator absorbs local latent noise into
plausible variations. (c) Residual fraction of the per-block Wiener
deprotector on held-out codes versus the analytic prediction; the two agree to
three decimals, so $\kappa_h$ is a measured quantity, not a fit artefact.}
\label{fig:latent}
\end{figure*}


% =====================================================================
\section{EXPERIMENTAL VALIDATION}
\label{sec:experiments}
% =====================================================================

\subsection{Implementation and Calibration}

The closed-form allocator, the four accounting solvers, the audit routines and
the unit tests accompany the paper; a one-command script reproduces every
number below. The KKT multiplier is found by Brent's root finder on $\log_{10}\nu\in[-18,18]$
(max.\ 200 iterations); the $(\varepsilon,\delta)$ conversion searches 140 RDP orders.
Three calibrations are used. The \emph{synthetic anchor} is anchored to the
reference backbone's clipped range and exercises every rule end to end; the
\emph{measured proxy} is built from 500 CelebAMask-HQ faces~\cite{celebamaskhq};
and the \emph{measured latent} calibration of Section~\ref{sec:latent} is built
on the released SemanticStyleGAN generator. The HF parquet
copy provides aligned 1024$\times$1024 images and 512$\times$512 part masks
(19 classes); the masks are upsampled nearest-neighbour to the image size, and
the 19 classes are grouped into the 15 blocks of Section~\ref{sec:setup}. Per
block we measure the within-part pixel variance $\bar g_h$ and the mean 0.5\%--
99.5\% clipped intensity range; the 0.5\%/99.5\% quantiles trim the extreme
pixels so the range is not set by a handful of outliers. Across the 500 images
the per-part pixel variance lies in $[0.025,0.066]$ and the clipped range in
$[0.67,0.90]$; the per-block values are listed in the accompanying
\texttt{part\_stats.csv}. The trace proxy is $g_h=d_h\bar g_h$, and the
sensitivity is $\Delta f_h=\sqrt{d_h}\,\mathrm{range}_h$.

For the image-proxy calibration, the remaining quantities are surrogates, not
measurements on a trained encoder. The utility weight $w_h$ and residual
fraction $\kappa_h$ are seeded defaults (seed~0): $w_h$ is scaled by an
identity weight and perturbed by an $8\%$ Gaussian clipped to $[0.8,1.2]$, and
$\kappa_h$ by a $5\%$ Gaussian clipped to $[0.85,1.10]$ and then to
$[0.2,0.98]$. The LPIPS distortion and the attack error $P_e$ are
model-surrogate predictions derived from these defaults; they are not encoder
latent measurements, and no multi-seed mean or variance is reported because
the surrogate curves are generated under the single seed~0 calibration. These
proxy values are superseded by the latent-level measurement of
Section~\ref{sec:latent}, where $w_h,\kappa_h$ are measured on the released
generator; only $P_e$ remains a surrogate there.

\emph{Reproducibility.} The pipeline scripts take the data parquet, sample
count $n$, target floor $\tau$, budget $\varepsilon$, $\delta$, gallery size
$K$, static count $k$, and output directory as inputs; they write the per-part
statistics, calibration metadata, allocation summaries, PGFPlots `.dat'
sources, and the validation figures. A separate latent pipeline loads the
released SemanticStyleGAN CelebAMask-HQ checkpoint, samples prior codes through
the mapping network, runs the per-block noise sweeps, and writes the measured
calibration of Section~\ref{sec:latent}; the generator's fused CUDA operators
fall back to their native CPU kernels, so the full latent calibration
reproduces on CPU, and no measured quantity is hard-coded. The one measurement
that still requires a GPU-scale environment is the identification attack:
inverting a labeled identity corpus through the generator and training the
recognition head to produce a measured $P_e$.

\subsection{Certificate Price versus Task Overlap}

Table~\ref{tab:synthetic} reports the five allocators on the synthetic anchor
in the task-decoupled regime ($\tau=0.90$, $K=10^{4}$, $\varepsilon=100$,
$\delta=10^{-5}$). The MI certificate meets the 0.33-bit leakage budget and
holds the 0.900 floor; the RDP, GDP and basic solvers trade the same budget
for much larger Bob distortion; the static $k=4$ index selection keeps Bob's
LPIPS low but leaks $5.5\times10^{3}$ bit and certifies no floor at all.

\begin{table}[t]
\centering\footnotesize
\setlength{\tabcolsep}{3.5pt}
\caption{Synthetic anchor, task-decoupled ($\tau{=}0.90$, $K{=}10^{4}$,
$\varepsilon{=}100$, $\delta{=}10^{-5}$). The source powers and sensitivities
are anchored to the reference backbone's clipped range; LPIPS, residual
fraction and attack error are surrogate predictions. ``leak.'' is the certified
leaked mutual information in bits; ``floor'' is the Fano identification-error
floor; the GDP column lists the total privacy-loss parameter and the certified
$\varepsilon$.}
\label{tab:synthetic}
\begin{tabular}{lcccc}
\toprule
scheme & LPIPS$^\dagger$ & leak. (bit) & floor & GDP ($\mu_{\mathrm{tot}},\varepsilon$) \\
\midrule
MI certificate & 0.240 & 0.329 & 0.900 & 70.1,\,2755 \\
RDP & 0.793 & 0.007 & 0.924 & 10.13,\,94.5 \\
GDP & 0.788 & 0.007 & 0.924 & 10.51,\,100.0 \\
$\ell_1$ basic & 0.853 & 0.006 & 0.924 & 8.44,\,71.6 \\
static $k{=}4$ & 0.146 & $5.5\times10^{3}$ & 0.000 & --- \\
\bottomrule
\end{tabular}
$^\dagger$surrogate.
\end{table}

Table~\ref{tab:measured} reports the measured-proxy calibration on two
scenarios. In the reconstruction scenario the private semantics are the task
(task overlap 1.000, $|P|=15$); in the task-decoupled scenario the private
blocks are task-null (overlap 0.013, $|P|=6$). Under the same 0.900 floor the
decoupled scenario incurs a smaller proxy distortion than reconstruction
(0.771 versus 0.862), so the price of a certificate is set by the
task--private overlap, not by the reported $\varepsilon$. Because $w_h,\kappa_h$
are surrogates, these LPIPS magnitudes are not measured distortions and should
not be quoted against a no-protection baseline; the reliable readings are the
floor and the relative overlap ordering. Figure~\ref{fig:certificate}
plots the same price--floor curve for both calibrations.

The overlap determines the certificate price qualitatively. When the private
semantics project onto the same subspace the task needs, the noise that
protects identity also sits on the task coordinates: Bob's distortion rises
unavoidably, because there is no way to obscure the private information without
degrading the task signal. When the private blocks are task-null, they carry
no information Bob needs, so they can be protected with little damage to the
released task output; the decoupled scenario therefore reaches the same
0.900 floor at a smaller proxy distortion. The measured overlap values
(1.000 versus 0.013) encode this geometry, and the ordering holds on both
calibrations.

\emph{Gallery-size ablation.} Figure~\ref{fig:imax} sweeps the gallery size
$K\in\{10^{3},10^{4},10^{6}\}$ against the feasible floor $\tau$. Both the
ceiling $\tau_{\max}(K)$ and the budget $I_{\max}(\tau,K)$ grow only
logarithmically in $K$: enlarging the template bank does not buy a strong
certified floor for free, and a high error floor still requires a very small
leakage budget.

\begin{table}[t]
\centering\footnotesize
\setlength{\tabcolsep}{3.5pt}
\caption{Measured proxy, 500 CelebAMask-HQ faces ($\tau{=}0.90$, $K{=}10^{4}$).
$g_h,\Delta f_h$ measured; $w_h,\kappa_h,P_e$ surrogate; LPIPS values are
proxy predictions, not measured distortions. ``leak.'' is the certified leaked
mutual information in bits; ``floor'' is the Fano identification-error floor.}
\label{tab:measured}
\begin{tabular}{lccccc}
\toprule
scenario & overlap & $|P|$ & MI LPIPS$^\dagger$ & leak. (bit) & floor \\
\midrule
reconstruction & 1.000 & 15 & 0.862 & 0.329 & 0.900 \\
task-decoupled & 0.013 & 6 & 0.771 & 0.329 & 0.900 \\
\bottomrule
\end{tabular}
$^\dagger$surrogate.
\end{table}

\subsection{Audit of the Accounting Laws}

Figure~\ref{fig:composite}(b) audits the three composition laws on the same
release over the noise-to-signal range. At the low-noise (large-$\mu$)
operating point certified at quadrature $\varepsilon=2000$, $\delta=10^{-5}$,
the synthetic anchor gives $\mu_h\approx24.1$, $|P|=6$: the amplitude-only
$\ell_1$ form reports $\approx692.6$, the rigorous tail bound reports
$\approx2440.5$, and quadrature sits at 2000. Table~\ref{tab:audit} reproduces
the ordering on the measured proxy, adding the order-optimised RDP conversion
as a third conservative bound: the amplitude-only law under-reports
$\varepsilon$ once $\mu_h\gg c$, the RDP conversion sits slightly above
quadrature, and the exact tail bound stays above quadrature in both scenarios.
The divergence is regime-dependent: at the high-noise (small-$\mu$) end the
amplitude-only truncation approaches the \emph{rigorous linear tail}
(both dominated by the linear $c\mu_h$ term), yet it remains larger than
quadrature by a factor of about $\sqrt{|P|}$ (equal-$\mu$ case); the large gap
reported above appears only in the low-noise, large-$\mu$ regime. The red band in
Figure~\ref{fig:composite}(b) marks this regime, which corresponds to the
low-noise operating point a wiretap scheme prefers for Bob's fidelity; it is
precisely this high-fidelity regime where the amplitude-only audit understates
the released loss.

\begin{table}[t]
\centering\footnotesize
\setlength{\tabcolsep}{3.0pt}
\caption{Audit at quadrature $\varepsilon{=}2000$, $\delta{=}10^{-5}$.
Columns: $\bar\mu_h$ is the mean per-block privacy-loss parameter;
``amplitude'' is the amplitude-only $\ell_1$ truncation; ``RDP'' is the
order-optimised RDP to $(\varepsilon,\delta)$ conversion; ``quadrature'' is the
exact Gaussian-DP account; ``exact tail'' keeps the full linear bound
$\sum_h(\mu_h^2/2+c\mu_h)$. The exact tail lies above the GDP quadrature
account because it retains the per-block linear $c\mu_h$ term and does not use
the exact Gaussian privacy-loss distribution; the tail also sums per-block
$\mu_h^2/2+c\mu_h$, so it cannot be reconstructed from the reported mean
$\bar\mu_h$ alone.}
\label{tab:audit}
\begin{tabular}{lccccc}
\toprule
scenario & $\bar\mu_h$ & amplitude & RDP & quadrature & exact tail \\
\midrule
reconstruction & 15.0 & 1089.2 & 2037.8 & 2000 & 2837.0 \\
task-decoupled & 24.1 & 701.1 & 2037.8 & 2000 & 2448.9 \\
\bottomrule
\end{tabular}
\end{table}

\subsection{Static Index Selection}

The static $k=4$ rule selects, by the score $\Delta f_h/a_h$ (high sensitivity
at low task cost), the four cheapest components to protect and leaves the rest
untouched. Because the corrected certificate bound of Lemma~\ref{lm:maxent}
accounts for all $d_h$ coordinates of each cleanly released block, the leaked
mutual information is enormous: $5.5\times10^{3}$ bit on the synthetic anchor,
$7.7\times10^{3}$ bit in the measured decoupled scenario and $7.8\times10^{4}$ bit
in reconstruction (all far above the $0.33$-bit budget), and the Fano floor is
zero in every case. The defect is structural: only the $k=4$
selected private blocks receive noise, while every other private block is
released clean and still carries its identity signal. The certificate cannot
observe these untouched blocks, so their leakage accumulates without any
constraint and the identification floor drops to zero. As
Corollary~\ref{co:static} predicts, the
static rule is a degenerate corner of the allocation program: it looks cheap
only because the certificate cannot see what it leaves unprotected. Earlier
reports in which such a static rule looks adequate judged reconstruction
distortion alone; without an identification-error floor, the cleanly released
private blocks are invisible to that metric, which masks the identity leak.

\subsection{Measured Quantities versus Surrogates}

The measured and surrogate quantities are separated explicitly. On the
image-proxy calibration, the block source power $g_h\approx d_h\bar g_h$
($\bar g_h$ the within-part pixel variance, $d_h$ the code width) and
$\Delta f_h$ (part clipped range) are computed from real masks, while $w_h$,
$\kappa_h$ and $P_e$ are default surrogates. The latent-level calibration of
Section~\ref{sec:latent} replaces all of $g_h,\Delta f_h,w_h,\kappa_h$ by
measurements on the released generator; the only quantity still awaiting a
measurement is the identification error $P_e$, which requires an
inverted-identity corpus with labels and a recognition head.

\subsection{Latent-Level Measurement on the SemanticStyleGAN Backbone}
\label{sec:latent}

\emph{Protocol.} The backbone is the released SemanticStyleGAN
CelebAMask-HQ $512{\times}512$ checkpoint~\cite{semanticstylegan,celebamaskhq}.
The decomposed code follows the paper layout: two shared codes of width
$d_h=512$ and thirteen local part codes of width $d_h=1024$ (28 latent slots
of 512 coordinates), with block semantics background, skin, eyes, brows,
mouth, nose, ears, hair, neck, cloth, eye\_g, hat, ear\_r as ordered by the
dataset preprocessing. The source statistics are exact prior statistics:
$65{,}536$ $w$-codes drawn from the latent prior are propagated through the
real mapping network, giving the per-block source power $g_h$ and the
0.5\%/99.5\% clipped sensitivity $\Delta f_h$ of the released distribution;
this is a pushforward of the generator prior, not an inversion of real
identities. The utility weights are measured by a per-block noise sweep: for
each block $h$ and each $\sigma_h^2$ on a six-level grid, noise is added to
$N=16$ codes in block $h$ only, the generator decodes back to images, and
LPIPS~\cite{lpips} against the clean decode is fitted to
$\mathrm{LPIPS}\approx\mathrm{lpips}_0+w_h\sigma_h^2$. The residual fraction
$\kappa_h$ is the fitted residual slope of a per-block linear-MMSE (Wiener)
deprotector whose covariance is measured on the same $65{,}536$ codes; the
empirical residual on 1024 held-out codes matches the analytic Wiener
prediction to three decimals. All runs execute on CPU, with the generator's
fused CUDA operators falling back to their native CPU kernels.

\emph{Measured scalars.} Table~\ref{tab:latent_scalars} lists the fifteen
measured blocks. The per-coordinate standard deviation is $\approx0.697$
uniformly across blocks ($g_h\approx248.6$--$249.4$ for $d_h=512$ and
$\approx497.6$--$498.8$ for $d_h=1024$), while the clipped per-coordinate
range is $\approx4.32$ against the Gaussian prediction $3.59$: the mapping
network's output is heavier-tailed than Gaussian, so the clipped sensitivity
($\Delta f_h\approx97.6$--$97.8$ at $d_h=512$ and $\approx138.0$--$138.3$ at
$d_h=1024$; global sensitivity $\approx517$) exceeds what a Gaussian-range
assumption would certify. The measured $\kappa_h\approx0.24$ is essentially
flat across blocks, as the Wiener residual depends mainly on the local
signal-to-noise geometry. Figure~\ref{fig:latent} plots the utility
weights, example sweep curves with their quadratic fits, and the empirical
versus analytic $\kappa$ curves.

\begin{table}[t]
\centering\footnotesize
\setlength{\tabcolsep}{3.5pt}
\caption{Measured per-block scalars on the released SemanticStyleGAN
generator. $g_h,\Delta f_h$: $65{,}536$ prior codes through the real mapping
network; $w_h$: decode-and-LPIPS noise sweep ($N{=}16$ codes per block, six
noise levels); $\kappa_h$: fitted Wiener residual slope. The image-space
proxy of Table~\ref{tab:measured} is retained as an offline fallback.}
\label{tab:latent_scalars}
\begin{tabular}{lccccc}
\toprule
block & $d_h$ & $g_h$ & $\Delta f_h$ & $w_h$ & $\kappa_h$ \\
\midrule
shared\_0  & 512  & 248.65 & 97.63  & 0.2713 & 0.242 \\
shared\_1  & 512  & 249.35 & 97.77  & 0.3420 & 0.241 \\
background & 1024 & 497.58 & 138.08 & 0.1702 & 0.242 \\
skin       & 1024 & 497.69 & 138.18 & 0.1741 & 0.241 \\
eyes       & 1024 & 497.63 & 138.08 & 0.0097 & 0.242 \\
brows      & 1024 & 498.56 & 138.19 & 0.0159 & 0.243 \\
mouth      & 1024 & 498.76 & 138.25 & 0.0112 & 0.241 \\
nose       & 1024 & 498.31 & 138.13 & 0.0088 & 0.242 \\
ears       & 1024 & 497.78 & 138.12 & 0.0170 & 0.242 \\
hair       & 1024 & 497.64 & 138.16 & 0.2753 & 0.241 \\
neck       & 1024 & 497.57 & 138.01 & 0.0281 & 0.241 \\
cloth      & 1024 & 498.50 & 138.27 & 0.0424 & 0.241 \\
eye\_g     & 1024 & 498.63 & 138.25 & 0.0117 & 0.243 \\
hat        & 1024 & 497.75 & 138.13 & 0.0276 & 0.241 \\
ear\_r     & 1024 & 497.79 & 138.15 & 0.0078 & 0.240 \\
\bottomrule
\end{tabular}
\end{table}

\emph{Allocation on the measured calibration.}
Table~\ref{tab:latent_alloc} reruns the five allocators on the measured
scalars ($\tau{=}0.90$, $K{=}10^{4}$, $\varepsilon{=}100$,
$\delta{=}10^{-5}$) in the two scenarios of Table~\ref{tab:measured}. The
identity-relevant private set (the two shared codes together with skin,
eyes, brows, mouth and nose, $|P|{=}7$) overlaps the LPIPS-bearing subspace
substantially: the measured task--private overlap is $0.590$, far above the
$0.013$ of the image proxy, because the perceptual utility concentrates on
the shared, hair, skin and background blocks that also carry identity. Under
an LPIPS task the certificate is therefore not cheap: Bob's predicted
distortion saturates near the model ceiling in both scenarios. The reliable
readings are the measured $w_h,\kappa_h$ themselves, the floors, and the
noise budget: moving from reconstruction to the decoupled private set cuts
the total injected noise-to-signal mass from $49.7$ to $11.0$ ($4.5\times$)
and raises the allocation sparsity from $0.123$ to $0.249$. The static
$k{=}4$ rule again fails structurally, leaking $9.7\times10^{4}$ bit in
reconstruction and $2.9\times10^{4}$ bit in the decoupled scenario with a zero
floor.

\begin{table}[t]
\centering\footnotesize
\setlength{\tabcolsep}{3.5pt}
\caption{Measured latent calibration ($\tau{=}0.90$, $K{=}10^{4}$,
$\varepsilon{=}100$, $\delta{=}10^{-5}$). All four scalars measured on the
released generator; LPIPS$^\dagger$ is the saturating prediction
$\mathrm{lpips}_0+0.75\,S/(0.75+S)$ driven by the measured sweep and hits its
$0.86$ ceiling at $\tau{=}0.9$ in both scenarios. ``leak.'' is the certified
leaked mutual information in bits; ``floor'' is the Fano identification-error
floor.}
\label{tab:latent_alloc}
\begin{tabular}{llccc}
\toprule
scenario & scheme & LPIPS$^\dagger$ & leak. (bit) & floor \\
\midrule
reconstruction & MI certificate & 0.862 & 0.329 & 0.900 \\
(overlap 1.000, & RDP            & 0.861 & 1.931 & 0.779 \\
$|P|{=}15$)    & GDP            & 0.861 & 2.077 & 0.768 \\
               & $\ell_1$ basic & 0.862 & 0.634 & 0.877 \\
               & static $k{=}4$ & 0.832 & $9.7\times10^{4}$ & 0.000 \\
\midrule
task-decoupled & MI certificate & 0.861 & 0.329 & 0.900 \\
(overlap 0.590, & RDP           & 0.857 & 1.931 & 0.779 \\
$|P|{=}7$)     & GDP            & 0.857 & 2.076 & 0.769 \\
               & $\ell_1$ basic & 0.860 & 1.367 & 0.822 \\
               & static $k{=}4$ & 0.832 & $2.9\times10^{4}$ & 0.000 \\
\bottomrule
\end{tabular}
$^\dagger$saturating surrogate scale; measured inputs.
\end{table}

\emph{Findings.} Three observations survive the move from proxy to
measurement. First, the measured utility weights are concentrated:
$w_h$ is large on shared\_1 ($0.342$), hair ($0.275$), shared\_0 ($0.271$),
skin ($0.174$) and background ($0.170$), while the facial-feature blocks are
nearly free to protect ($w_{\mathrm{eyes}}{=}0.0097$,
$w_{\mathrm{nose}}{=}0.0088$, $w_{\mathrm{ear\_r}}{=}0.0078$). Visual
inspection of the sweep grids shows why: the generator absorbs local latent
noise into plausible variations, such as make-up-like changes around the eyes
at $\sigma_h{=}0.8$ or recolouring for hair, rather than perceptual distortion,
so a decomposed latent is not uniformly sensitive to noise. Second, identity
is entangled with utility through the shared and skin blocks, which is what
keeps the decoupled-scenario overlap at $0.590$; the certificate price is
again governed by overlap, not by $\varepsilon$. Third, the audit ordering of
Proposition~\ref{pr:quadrature} reproduces on the measured scalars: at
quadrature $\varepsilon{=}2000$, $\delta{=}10^{-5}$, the reconstruction
scenario ($\bar\mu_h\approx14.5$) reports amplitude-only $1055.2$ $<
$quadrature $2000$ $<$ RDP $2037.8$ $<$ exact tail $2803.0$, and the decoupled
scenario ($\bar\mu_h\approx21.2$) reports $718.4 < 2000 < 2037.8 < 2466.3$.

\emph{What remains pending.} The attack error $P_e$ is still the Fano-floor certificate evaluated against a
surrogate adversary. Measuring it would require inverting a labeled identity corpus
through the generator --- CelebAMask-HQ file names joined back to the CelebA identity
list --- and training a recognition head on held-out identities. That is a GPU-scale
inversion workload, which this CPU-only build does not attempt. The codes sampled from
the generator carry no real identity, so the released calibration data raises no
subject-privacy concern.

% =====================================================================
\section{DISCUSSION AND LIMITATIONS}
\label{sec:discussion}
% =====================================================================

\subsection{Limitations}

The study carries six limitations. L1 and L6 are inherent to the certificate
framework; L2--L5 are limitations of the current implementation and calibration.

\begin{enumerate}
\item[(L1)] \emph{Attack scope (framework-inherent).} The certificate of
Theorem~\ref{th:certificate} is an identification guarantee, not a guarantee on
arbitrary attribute inference. It constrains every estimator of the template
$B$ but says nothing about an attribute that is not a function of $z_P$.
\item[(L2)] \emph{Surrogate attack error and prior-pushforward calibration
(implementation-level).}
The four calibration scalars are measured on the released generator
(Section~\ref{sec:latent}), but the identification error $P_e$ is still a
surrogate evaluated against the Fano floor, not a trained recognition head on
held-out identities. In addition, the measured source statistics come from
prior codes pushed through the real mapping network, not from codes inverted
from real identities; an inversion-based calibration could shift $g_h$ and
$\Delta f_h$ if real codes lie off the prior manifold.
\item[(L3)] \emph{Manual private partition (implementation-level).} The private subspace $P$ is
declared by the operator before the noise is drawn. The method does not address
automatic learning or partitioning of the private semantic components.
\item[(L4)] \emph{Conservative correlation bound (implementation-level).} The block-wise second-order
structure of Assumption~\ref{as:indep} ignores residual within- and
across-block correlation; the AM--GM second-order bound is a safe upper bound but makes the
estimated utility performance look pessimistic (Remark~\ref{rm:corr}).
\item[(L5)] \emph{Key-conditional mode (implementation-level).} When the shaping $T$ is a Bob-secret
key unknown to the eavesdropper, the resulting guarantee is key-conditional
privacy, not plain $(\varepsilon,\delta)$-DP under the Kerckhoffs assumption
(Remark~\ref{rm:tension}).
\item[(L6)] \emph{Finite template bank (framework-inherent).} For any finite $K$, the feasibility
ceiling $\tau_{\max}(K)<1$; an arbitrarily strong certified identification
error upper bound is unattainable, because the template entropy itself is the
budget currency.
\end{enumerate}

\subsection{Future Work}

Four directions follow directly from these limitations.

\begin{enumerate}
\item[(F1)] \emph{Measured identification error on inverted identities.} Two of the three latent-level validation stages are complete
(Section~\ref{sec:latent}): the released SemanticStyleGAN generator fixes the real per-block
variance, clipped range, utility weights and Wiener residuals, and checks
Assumptions~\ref{as:range}--\ref{as:indep} on the true latent geometry. The attack stage remains. Invert a
labeled identity corpus (CelebAMask-HQ file names joined to the CelebA identity list)
through the generator, train a recognition head on held-out identities, and measure the
identification error against the certified floor, replacing the surrogate $P_e$ used here.
This would also upgrade the source statistics from a prior pushforward to an
inversion-based calibration. \item[(F2)] \emph{Automatic private-partition learning.} Develop a method that derives the private
set $P$ from an explicit privacy target and the measured task--private overlap, instead
of declaring it by hand. \item[(F3)] \emph{Accounting beyond Gaussian noise.} Extend the auditable composition and
certificate to non-Gaussian noise mechanisms, where the privacy-loss distribution is not
exactly Gaussian and GDP quadrature no longer applies. \item[(F4)] \emph{Tightening the correlated-source bound.} Under strong
between-block correlation, estimate the joint covariance and replace the block-wise
AM--GM relaxation by the exact determinant bound of Lemma~\ref{lm:maxent}, so the certified leakage
reflects the true joint mutual information more tightly. \item[(F5)] \emph{Cross-release composition.} Study budget
composition over multiple releases to the same subject; the present analysis is
single-release and does not cover a repeated-query adversary.
\end{enumerate}

% =====================================================================
\section{CONCLUSION}
\label{sec:conclusion}
% =====================================================================

Three practical conclusions follow from this study. First, the privacy of a
semantic link cannot be judged from the appearance of the released noise; an
auditable budget must be computed from the component sensitivity and the
composition law, and reported alongside an adversary identification bound
(the two guarantees are independent). Second, amplitude-only $\ell_1$
accounting inherits the Laplace-mechanism form and under-reports the leakage of
a Gaussian release in the low-noise (large-$\mu$) regime, presenting a
tighter-looking assessment than the mechanism guarantees. Third, a dynamic,
globally coordinated budget allocation outperforms protecting a fixed subset of
components: the static rule leaves the untouched private blocks released clean
and certifies no identification floor. The framework is a mechanism-level,
general-purpose tool for auditing latent releases that satisfy the stated
model assumptions, and is not tied to SemanticStyleGAN; that generator serves
only as the testbed for the numerical calibration.

The evidence now includes a latent-level measurement on the released SemanticStyleGAN
generator. All four calibration scalars are measured on real decomposed codes and
generator decodes. Three results reproduce on that measured calibration: the accounting
ordering, the overlap--price relation, and the static-rule failure. One surrogate
remains --- the identification attack error. Measuring it on an inverted-identity corpus
is the first item of future work.

\section*{CODE AND DATA AVAILABILITY}
The calibration routines, the closed-form allocator and its accounting solvers, the
reproduction pipeline, the latent-level measurement pipeline, the PGFPlots data files
and the unit tests are available at
\href{https://github.com/wqinmin975/Auditable-Privacy-Accounting-and-Dynamic-Semantic-Budget}{\texttt{github.com/wqinmin975/\allowbreak Auditable-\allowbreak Privacy-\allowbreak Accounting-\allowbreak and-\allowbreak Dynamic-\allowbreak Semantic-\allowbreak Budget}};
the archived release DOI will be inserted upon publication. The CelebAMask-HQ shard used
for the measured proxy is downloaded from the public Hugging Face dataset \texttt{v-xchen-v/celebamask\_hq}; it is
used under the original dataset's research license, for proxy statistics only. The
latent calibration loads the official SemanticStyleGAN CelebAMask-HQ checkpoint from the
authors' public release, with the generator code vendored under \texttt{third\_party/} together with its
license. No face images, identity labels, or inverted latent corpus are redistributed
here. The reported latent statistics come from generator-sampled prior codes that carry
no real identity, and only aggregate per-block scalars and curves are reported.
Reproduction splits in two. The synthetic anchor and proxy calibrations run entirely on
CPU from the committed code and `.dat` sources. The full identification-attack
experiment needs a GPU and an inverted, labeled face-identity corpus, and is not yet
implemented in the release. Large generated binaries --- multi-GB source-statistics
arrays and the pretrained checkpoint --- are excluded by \texttt{.gitignore}; the repository README
documents how to regenerate or download each of them.

\begin{thebibliography}{00}\leftskip1pc
\bibitem{dong2022gdp} J.~Dong, A.~Roth, and W.~J. Su,
``Gaussian differential privacy,'' \emph{J. Roy. Stat. Soc. B}, vol.~84,
no.~1, pp.~3--37, 2022.
\bibitem{mironov2017rdp} I.~Mironov, ``R\'enyi differential privacy,'' in
\emph{Proc. IEEE 30th Comput. Secur. Found. Symp. (CSF)}, 2017, pp.~263--275.
\bibitem{dwork2014} C.~Dwork and A.~Roth, \emph{The Algorithmic Foundations of
Differential Privacy}. Hanover, MA, USA: Now Publishers, 2014.
\bibitem{abadi2016} M.~Abadi \emph{et al.}, ``Deep learning with differential
privacy,'' in \emph{Proc. ACM SIGSAC Conf. Comput. Commun. Secur. (CCS)}, 2016,
pp.~308--318.
\bibitem{shokri2017} R.~Shokri, M.~Stronati, C.~Song, and V.~Shmatikov,
``Membership inference attacks against machine learning models,'' in
\emph{Proc. IEEE Symp. Secur. Privacy (S\&P)}, 2017, pp.~3--18.
\bibitem{fano} T.~M. Cover and J.~A. Thomas, \emph{Elements of Information
Theory}, 2nd ed. Hoboken, NJ, USA: Wiley, 2006.
\bibitem{privacyfunnel} A.~Makhdoumi, S.~Salamatian, N.~Fawaz, and M.~M\'edard,
``From the information bottleneck to the privacy funnel,'' in \emph{Proc. IEEE
Inf. Theory Workshop (ITW)}, 2014, pp.~501--505.
\bibitem{semcomsurvey_won} D.~Won, G.~Woraphonbenjakul, A.~B.~Wondmagegn,
A.-T.~Tran, D.~Lee, D.~S.~Lakew, and S.~Cho, ``Resource management, security,
and privacy issues in semantic communications: A survey,'' \emph{IEEE Commun.
Surveys Tutorials}, vol.~27, no.~3, pp.~1758--1797, 2025.
\bibitem{semcomsurvey_meng} R.~Meng \emph{et al.}, ``A survey of secure semantic
communications,'' \emph{J. Netw. Comput. Appl.}, vol.~239, Art. no.~104181,
2025, doi: 10.1016/j.jnca.2025.104181.
\bibitem{semanticstylegan} Y.~Shi, X.~Yang, Y.~Wan, and X.~Shen,
``SemanticStyleGAN: Learning compositional generative priors for controllable
image synthesis and editing,'' in \emph{Proc. IEEE/CVF Conf. Comput. Vis.
Pattern Recognit. (CVPR)}, 2022, pp.~11254--11264.
% Title, authors, Δf and ε grid verified against arXiv:2510.23274v2 (2026-04-23).
\bibitem{arXiv2510} W.~Chen, Q.~Yang, S.~Shao, S.~Tang, Z.~Shi, and S.~Yu,
``Privacy-preserving semantic communication over wiretap channels with learnable
differential privacy,'' 2025, revised 2026, \emph{arXiv:2510.23274v2}.
\bibitem{celebamaskhq} C.-H. Lee, Z.~Liu, L.~Wu, and P.~Luo, ``MaskGAN:
Towards diverse and interactive facial image manipulation,'' in \emph{Proc.
IEEE/CVF Conf. Comput. Vis. Pattern Recognit. (CVPR)}, 2020, pp.~5549--5558.
\bibitem{lpips} R.~Zhang, P.~Isola, A.~A. Efros, E.~Shechtman, and O.~Wang,
``The unreasonable effectiveness of deep features as a perceptual metric,'' in
\emph{Proc. IEEE/CVF Conf. Comput. Vis. Pattern Recognit. (CVPR)}, 2018,
pp.~586--595.
\end{thebibliography}

\begin{IEEEbiographynophoto}
{Qinmin Wang} received the M.Sc. degree from the National Digital Switching
System Engineering and Technological R\&D Center (NDSC) in 2006 and the Ph.D.
degree from the University of Electronic Science and Technology of China
(UESTC) in 2009. He is currently an Associate Professor with the School of
Computer Science and Artificial Intelligence, Guangdong University of Science
and Technology, Dongguan, China. His research interests include semantic
communications, communication anti-jamming techniques, radar signal
processing, non-cooperative communications, and deep learning.
\end{IEEEbiographynophoto}

\begin{IEEEbiographynophoto}
{Zixin Wang} is with the College of Electronic and Communication Engineering,
Tianjin Normal University, Tianjin, China. His research interests include
semantic communications and communication signal processing.
\end{IEEEbiographynophoto}

\begin{IEEEbiographynophoto}
{Jie Peng} (born 1981) received the M.Sc. degree. He is with Shenzhen Customs, Shenzhen,
China.
\end{IEEEbiographynophoto}

\begin{IEEEbiographynophoto}
{Fenggui Pei} is with the School of Computer Science and Artificial
Intelligence, Guangdong University of Science and Technology, Dongguan, China.
\end{IEEEbiographynophoto}

\begin{IEEEbiographynophoto}
{Zongzhong Wang} (Corresponding author) is with the School of
Electromechanical Engineering, Guangdong University of Technology, Guangzhou,
China.
\end{IEEEbiographynophoto}


\end{document}
